

%%%%%%%%%%%%%%%%%%%%%%%%%%%%%%%%%%%%%%%%%%%%%%%%%%%%%%%%%%%%%%%
% Analysis 1
%%%%%%%%%%%%%%%%%%%%%%%%%%%%%%%%%%%%%%%%%%%%%%%%%%%%%%%%%%%%%%%


\begin{enumerate}
  \item Prove the following properties of \(\setR\) from the definition of \((\setR,+)\) and \((\setR^*,\cdot)\) as abelian groups where multiplication distributes over addition: \(a(b+c)=ab+ac\) for all \(a,b,c\in\setR\).
  \begin{enumerate}
    \item \(a\cdot 0=0\).
    \begin{subprove}
      Let \(a \in \setR\).
      It will be proved that \(a \cdot 0 = 0\).
      
      Since \(0 + 0 = 0\), by the distributive property, \(a \cdot (0 + 0) = a \cdot 0 + a \cdot 0\).
      Therefore, \(a \cdot 0 = a \cdot 0 + a \cdot 0\).
      By adding the additive inverse of \(a \cdot 0\) to both sides, it follows that \(0 = a \cdot 0\).
      
      Therefore, \(a \cdot 0 = 0\).
    \end{subprove}
    \item \((-a)b=-ab\).
    \begin{subprove}
      Let \(a, b \in \setR\).
      It will be proved that \((-a) b = - (a b)\).
      
      Since \(0 = 0 \cdot b\), by the previous result, \(0 = (a + (-a)) b = a b + (-a) b\).
      By adding the additive inverse of \(a b\) to both sides, it follows that \((-a) b = - (a b)\).
      
      Therefore, \((-a) b = - (a b)\).
    \end{subprove}
    \item \((-a)(-b)=ab\).
    \begin{subprove}
      Let \(a, b \in \setR\).
      It will be proved that \((-a)(-b) = a b\).
      
      Since \(0 = 0 \cdot (-b)\), by the previous result, \(0 = (a + (-a))(-b) = a(-b) + (-a)(-b)\).
      By adding the additive inverse of \(a(-b)\) to both sides, it follows that \((-a)(-b) = - (a(-b))\).
      By the previous result again, \(- (a(-b)) = - (-(a b)) = a b\).
      
      Therefore, \((-a)(-b) = a b\).
    \end{subprove}
    \item If \(ab=0\), then either \(a=0\) or \(b=0\).
    \begin{subprove}
      Let \(a, b \in \setR\) such that \(a b = 0\).
      It will be proved that either \(a = 0\) or \(b = 0\).
      
      Suppose that \(a \neq 0\).
      Since \((\setR^*, \cdot)\) is a group, \(a\) has a multiplicative inverse \(a^{-1} \in \setR^*\).
      By multiplying both sides of the equation \(a b = 0\) by \(a^{-1}\), it follows that \(b = a^{-1} \cdot 0 = 0\).
      
      Therefore, if \(a b = 0\) and \(a \neq 0\), then \(b = 0\).
      Hence, if \(a b = 0\), then either \(a = 0\) or \(b = 0\).
    \end{subprove}
    \item \(-a=(-1)a\).
    \begin{subprove}
      Let \(a \in \setR\).
      It will be proved that \(-a = (-1) a\).
      
      Since \(0 = 0 \cdot a\), by the previous result, \(0 = (1 + (-1)) a = 1 \cdot a + (-1) a = a + (-1) a\).
      By adding the additive inverse of \(a\) to both sides, it follows that \(-a = (-1) a\).
      
      Therefore, \(-a = (-1) a\).
    \end{subprove}
  \end{enumerate}

  \item Prove the following properties of \(\setR\) from the group definitions above, together with the total order \(\le\) on \(\setR\) satisfying: for all \(a,b,c\in\setR\), if \(a\le b\) then \(a+c\le b+c\); and if \(a\le b\) and \(0\le c\) then \(ac\le bc\).
  \begin{enumerate}
    \item If \(a\le b\), then \(-b\le -a\).
    \begin{subprove}
      Let \(a, b \in \setR\) such that \(a \le b\).
      It will be proved that \(-b \le -a\).
      
      Since \(a \le b\), by adding \(-b\) to both sides, it follows that \(a + (-b) \le b + (-b)\).
      Therefore, \(a - b \le 0\).
      By adding \(-a\) to both sides, it follows that \(-b \le -a\).
      
      Therefore, if \(a \le b\), then \(-b \le -a\).
    \end{subprove}
    \item If \(a\le b\) and \(c\le 0\), then \(bc\le ac\).
    \begin{subprove}
      Let \(a, b, c \in \setR\) such that \(a \le b\) and \(c \le 0\).
      It will be proved that \(b c \le a c\).
      
      Since \(c \le 0\), by multiplying both sides of the inequality \(a \le b\) by \(c\) (which is nonpositive), the direction of the inequality reverses.
      Therefore, \(b c \le a c\).
      
      Therefore, if \(a \le b\) and \(c \le 0\), then \(b c \le a c\).
    \end{subprove}
    \item If \(0\le a\) and \(0\le b\), then \(0\le ab\).
    \begin{subprove}
      Let \(a, b \in \setR\) such that \(0 \le a\) and \(0 \le b\).
      It will be proved that \(0 \le a b\).
      
      Since \(0 \le a\), by multiplying both sides of the inequality \(0 \le b\) by \(a\) (which is nonnegative), it follows that \(0 \le a b\).
      
      Therefore, if \(0 \le a\) and \(0 \le b\), then \(0 \le a b\).
    \end{subprove}
    \item \(0\le a^2\).
    \begin{subprove}
      Let \(a \in \setR\).
      It will be proved that \(0 \le a^{2}\).
      
      By the trichotomy property of the total order on \(\setR\), either \(a < 0\), \(a = 0\), or \(a > 0\).
      If \(a = 0\), then \(a^{2} = 0^{2} = 0\), so \(0 \le a^{2}\).
      If \(a > 0\), then by the previous result, \(0 \le a^{2}\).
      If \(a < 0\), then \(-a > 0\), so by the previous result, \(0 \le (-a)^{2} = a^{2}\).
      
      Therefore, in all cases, \(0 \le a^{2}\).
    \end{subprove}
    \item \(0<1\).
    \begin{subprove}
      It will be proved that \(0 < 1\).
      
      Suppose, for the sake of contradiction, that \(1 \le 0\).
      By multiplying both sides of the inequality \(1 \le 0\) by \(-1\) (which is negative), the direction of the inequality reverses.
      Therefore, \(-1 \ge 0\).
      By adding \(1\) to both sides, it follows that \(0 \ge 1\).
      This contradicts the assumption that \(1 \le 0\).
      
      Therefore, it must be that \(0 < 1\).
    \end{subprove}
    \item If \(0<a\), then \(0<a^{-1}\).
    \begin{subprove}
      Let \(a \in \setR\) such that \(0 < a\).
      It will be proved that \(0 < a^{-1}\).
      
      Suppose, for the sake of contradiction, that \(a^{-1} \le 0\).
      By multiplying both sides of the inequality \(a^{-1} \le 0\) by \(a\) (which is positive), it follows that \(1 \le 0\).
      This contradicts the previous result that \(0 < 1\).
      
      Therefore, it must be that \(0 < a^{-1}\).
    \end{subprove}
    \item If \(0<a<b\), then \(0<b^{-1}<a^{-1}\).
    \begin{subprove}
      Let \(a, b \in \setR\) such that \(0 < a < b\).
      It will be proved that \(0 < b^{-1} < a^{-1}\).
      
      By the previous result, since \(0 < b\), it follows that \(0 < b^{-1}\).
      Since \(a < b\), by multiplying both sides of the inequality \(a < b\) by \(b^{-1}\) (which is positive), it follows that \(a b^{-1} < 1\).
      By multiplying both sides of the inequality \(a b^{-1} < 1\) by \(a^{-1}\) (which is positive), it follows that \(b^{-1} < a^{-1}\).
      
      Therefore, if \(0 < a < b\), then \(0 < b^{-1} < a^{-1}\).
    \end{subprove}
  \end{enumerate}

  \item Let \(a\in\setR\). The absolute value \(|a|\) is defined by \(|a|=a\) if \(a\ge0\) and \(|a|=-a\) if \(a<0\). Explain why the following follow directly from this definition.
  \begin{enumerate}
    \item \(|a|\ge0\) and \(|a|=0\) iff \(a=0\).
    \begin{subprove}
      Let \(a \in \setR\).
      It will be proved that \(|a| \ge 0\) and \(|a| = 0\) if and only if \(a = 0\).
      
      By the definition of absolute value, if \(a \ge 0\), then \(|a| = a \ge 0\).
      If \(a < 0\), then \(|a| = -a > 0\).
      Therefore, in both cases, \(|a| \ge 0\).
      
      Furthermore, if \(a = 0\), then \(|a| = |0| = 0\).
      Conversely, if \(|a| = 0\), then by the definition of absolute value, either \(a = 0\) (if \(a \ge 0\)) or \(-a = 0\) (if \(a < 0\)), which also implies \(a = 0\).
      
      Therefore, \(|a| \ge 0\) and \(|a| = 0\) if and only if \(a = 0\).
    \end{subprove}
    \item For \(a,b\in\setR\), \(|ab|=|a|\cdot|b|\).
    \begin{subprove}
      Let \(a, b \in \setR\).
      It will be proved that \(|ab| = |a| \cdot |b|\).
      
      By the definition of absolute value, the signs of \(a\) and \(b\) determine the values of \(|a|\) and \(|b|\).
      The following cases are considered.
      
      \textbf{Case 1:} Suppose \(a \ge 0\) and \(b \ge 0\).
      Then, by the definition of absolute value, \(|a| = a\) and \(|b| = b\).
      Since \(a \ge 0\) and \(b \ge 0\), it follows that \(ab \ge 0\).
      Therefore, \(|ab| = ab = |a| \cdot |b|\).
      
      \textbf{Case 2:} Suppose \(a \ge 0\) and \(b < 0\).
      Then, by the definition of absolute value, \(|a| = a\) and \(|b| = -b\).
      Since \(a \ge 0\) and \(b < 0\), it follows that \(ab < 0\).
      Therefore, \(|ab| = -(ab) = -ab = a(-b) = |a| \cdot |b|\).
      
      \textbf{Case 3:} Suppose \(a < 0\) and \(b \ge 0\).
      Then, by the definition of absolute value, \(|a| = -a\) and \(|b| = b\).
      Since \(a < 0\) and \(b \ge 0\), it follows that \(ab < 0\).
      Therefore, \(|ab| = -(ab) = -ab = (-a)b = |a| \cdot |b|\).
      
      \textbf{Case 4:} Suppose \(a < 0\) and \(b < 0\).
      Then, by the definition of absolute value, \(|a| = -a\) and \(|b| = -b\).
      Since \(a < 0\) and \(b < 0\), it follows that \(ab > 0\).
      Therefore, \(|ab| = ab = (-a)(-b) = |a| \cdot |b|\).
      
      \textbf{Case 5:} Suppose \(a = 0\) or \(b = 0\).
      Then \(ab = 0\), so \(|ab| = 0\).
      If \(a = 0\), then \(|a| = 0\), so \(|a| \cdot |b| = 0 \cdot |b| = 0\).
      If \(b = 0\), then \(|b| = 0\), so \(|a| \cdot |b| = |a| \cdot 0 = 0\).
      Therefore, \(|ab| = 0 = |a| \cdot |b|\).
      
      Therefore, in all cases, \(|ab| = |a| \cdot |b|\).
    \end{subprove}
    \item For \(a\in\setR\), \(|a|^2=a^2\).
    \begin{subprove}
      Let \(a \in \setR\).
      It will be proved that \(|a|^2 = a^2\).
      
      By the definition of absolute value, two cases are considered.
      
      \textbf{Case 1:} Suppose \(a \ge 0\).
      Then, by the definition of absolute value, \(|a| = a\).
      Therefore, \(|a|^2 = a^2\).
      
      \textbf{Case 2:} Suppose \(a < 0\).
      Then, by the definition of absolute value, \(|a| = -a\).
      Therefore, \(|a|^2 = (-a)^2 = a^2\).
      
      Therefore, in both cases, \(|a|^2 = a^2\).
    \end{subprove}
  \end{enumerate}

  \item Prove the following properties of the absolute value.
  \begin{define}{Triangle Inequality}[def:triangle-inequality][red]
    For all \(a,b\in\setR\), \(|a+b|\le|a|+|b|\).
    This inequality is called the \textbf{triangle inequality}.
  \end{define}
  \begin{enumerate}
    \item For all \(a,b\in\setR\), \(|a+b|\le|a|+|b|\).
    \begin{subprove}
      Let \(a, b \in \setR\).
      It will be proved that \(|a+b| \le |a| + |b|\).
      
      By the definition of absolute value, for any \(x \in \setR\), we have \(-|x| \le x \le |x|\).
      In particular, \(-|a| \le a \le |a|\) and \(-|b| \le b \le |b|\).
      Adding these inequalities, it follows that
      \[
        -(|a| + |b|) = -|a| - |b| \le a + b \le |a| + |b|.
      \]
      Therefore, in both cases, \(|a + b| \le |a| + |b|\).
    \end{subprove}
    \item For all \(a,b\in\setR\), \(|b|\le a\) iff \(-a\le b\le a\).
    \begin{subprove}
      Let \(a, b \in \setR\).
      It will be proved that \(|b| \le a\) if and only if \(-a \le b \le a\).
      
      Note that if \(|b| \le a\), then \(a \ge 0\) since \(|b| \ge 0\).
      
      (\(\Rightarrow\)) Suppose that \(|b| \le a\) for the sake of direct proof.
      Since \(|b| \ge 0\) and \(|b| \le a\), it follows that \(a \ge 0\).
      By the definition of absolute value, two cases are considered.
      
      \textbf{Case 1:} Suppose \(b \ge 0\).
      Then, by the definition of absolute value, \(|b| = b\).
      Since \(|b| \le a\), it follows that \(b \le a\).
      Since \(b \ge 0\) and \(a \ge 0\), it follows that \(-a \le 0 \le b \le a\).
      Therefore, \(-a \le b \le a\).
      
      \textbf{Case 2:} Suppose \(b < 0\).
      Then, by the definition of absolute value, \(|b| = -b\).
      Since \(|b| \le a\), it follows that \(-b \le a\), which implies \(b \ge -a\).
      Since \(b < 0\) and \(a \ge 0\), it follows that \(-a \le b < 0 \le a\).
      Therefore, \(-a \le b \le a\).
      
      Therefore, if \(|b| \le a\), then \(-a \le b \le a\).
      
      (\(\Leftarrow\)) Suppose that \(-a \le b \le a\) for the sake of direct proof.
      Since \(-a \le b \le a\), it follows that \(a \ge -a\), which implies \(a \ge 0\).
      By the definition of absolute value, two cases are considered.
      
      \textbf{Case 1:} Suppose \(b \ge 0\).
      Then, by the definition of absolute value, \(|b| = b\).
      Since \(b \le a\), it follows that \(|b| = b \le a\).
      
      \textbf{Case 2:} Suppose \(b < 0\).
      Then, by the definition of absolute value, \(|b| = -b\).
      Since \(-a \le b\), it follows that \(-b \le a\).
      Therefore, \(|b| = -b \le a\).
      
      Therefore, if \(-a \le b \le a\), then \(|b| \le a\).
      
      Therefore, \(|b| \le a\) if and only if \(-a \le b \le a\).
    \end{subprove}
    \item For all \(a,b\in\setR\), \(||a|-|b||\le|a-b|\).
    \begin{subprove}
      Let \(a, b \in \setR\).
      It will be proved that \(||a| - |b|| \le |a - b|\).
      
      By the \nameref{def:triangle-inequality}, \(|a| = |(a - b) + b| \le |a - b| + |b|\).
      It follows that \(|a| - |b| \le |a - b|\).
      
      Similarly, by the \nameref{def:triangle-inequality}, \(|b| = |(b - a) + a| \le |b - a| + |a| = |a - b| + |a|\).
      It follows that \(|b| - |a| \le |a - b|\).
      Therefore, \(-(|a| - |b|) = |b| - |a| \le |a - b|\).
      
      Thus, \(-|a - b| \le |a| - |b| \le |a - b|\).
      By the previous result, since \(-|a - b| \le |a| - |b| \le |a - b|\), 
      it follows that \(||a| - |b|| \le |a - b|\).
      
      Therefore, for all \(a, b \in \setR\), \(||a| - |b|| \le |a - b|\).
    \end{subprove}
    \item For \(a,x,\varepsilon\in\setR\), \(|x-a|<\varepsilon\) iff \(a-\varepsilon<x<a+\varepsilon\).
    \begin{subprove}
      Let \(a, x, \varepsilon \in \setR\).
      It will be proved that \(|x-a| < \varepsilon\) if and only if \(a-\varepsilon < x < a+\varepsilon\).
      
      (\(\Rightarrow\)) Suppose that \(|x-a| < \varepsilon\) for the sake of direct proof.
      By the definition of absolute value, \(|x-a| < \varepsilon\) means that either \(x-a < \varepsilon\) and \(x-a \ge 0\), or \(-(x-a) < \varepsilon\) and \(x-a < 0\).
      
      \textbf{Case 1:} Suppose \(x-a \ge 0\).
      Then \(x-a < \varepsilon\), so \(x < a+\varepsilon\).
      Since \(x-a \ge 0\), it follows that \(x \ge a > a-\varepsilon\).
      Therefore, \(a-\varepsilon < x < a+\varepsilon\).
      
      \textbf{Case 2:} Suppose \(x-a < 0\).
      Then \(-(x-a) = a-x < \varepsilon\), so \(a-\varepsilon < x\).
      Since \(x-a < 0\), it follows that \(x < a < a+\varepsilon\).
      Therefore, \(a-\varepsilon < x < a+\varepsilon\).
      
      Therefore, in both cases, if \(|x-a| < \varepsilon\), then \(a-\varepsilon < x < a+\varepsilon\).
      
      (\(\Leftarrow\)) Suppose that \(a-\varepsilon < x < a+\varepsilon\) for the sake of direct proof.
      It follows that \(-\varepsilon < x-a < \varepsilon\).
      By the definition of absolute value, since \(-\varepsilon < x-a < \varepsilon\), it follows that \(|x-a| < \varepsilon\).
      
      Therefore, \(|x-a| < \varepsilon\) if and only if \(a-\varepsilon < x < a+\varepsilon\).
    \end{subprove}
    \item If \(a,b,x,y\in\setR\) and \(a<x<b,\ a<y<b\), then \(|y-x|<b-a\).
    \begin{subprove}
      Let \(a, b, x, y \in \setR\) such that \(a < x < b\) and \(a < y < b\).
      It will be proved that \(|y-x| < b-a\).
      
      By the definition of absolute value, \(|y-x| = y-x\) if \(y \ge x\), and \(|y-x| = x-y\) if \(y < x\).
      
      \textbf{Case 1:} Suppose \(y \ge x\).
      Then \(|y-x| = y-x\).
      Since \(y < b\) and \(x > a\), it follows that \(y-x < b-a\).
      Therefore, \(|y-x| = y-x < b-a\).
      
      \textbf{Case 2:} Suppose \(y < x\).
      Then \(|y-x| = x-y\).
      Since \(x < b\) and \(y > a\), it follows that \(x-y < b-a\).
      Therefore, \(|y-x| = x-y < b-a\).
      
      Therefore, in both cases, \(|y-x| < b-a\).
    \end{subprove}
  \end{enumerate}



% In the standard construction of the real numbers the following completeness properties are equivalent:
% \begin{enumerate}
%   \item the least upper bound property (Completeness Axiom),
%   \item the greatest lower bound property,
%   \item Dedekind-type completeness, and
%   \item Cauchy completeness.
% \end{enumerate}
% That is, any one of these statements implies the others when working in \(\setR\).

% \begin{define}{Dedekind Completeness}[def:dedekind-completeness][red]
%   $\setR$ is \textbf{Dedekind complete}, that is, it satisfies the following \textbf{Dedekind completeness property}.\proofstep
%   Let $A,B\subseteq\setR$ be nonempty such that $a\le b$ for all $a\in A$ and $b\in B$.
%   Then there exists an element $c\in\setR$ such that $a\le c\le b$ for all $a\in A$ and $b\in B$.
% \end{define}


% \begin{define}{Cauchy Completeness}[def:cauchy-completeness][red]
%   $\setR$ is \textbf{Cauchy complete}, that is, it satisfies the following \textbf{Cauchy completeness property}.\proofstep
%   Every \nameref{def:cauchy-sequence} in $\setR$ converges to a limit in $\setR$.
% \end{define}


  \item Prove the following using the completeness axiom and the order of \(\setR\).
  \begin{define}{Completeness Axiom}[def:completeness-axiom][red]
    $\setR$ is \textbf{complete}, that is, it satisfies the following \textbf{least upper bound property}.\proofstep
    Let $S\subseteq\setR$ be nonempty and bounded above, 
    then there exists an element $t\in\setR$ such that
    \begin{enumerate}
      \item $s\le t$ for all $s \in S$, and
      \item $t\le u$ for all any upper bound of $S$, $u\in\setR$, that is, $s\le u$ for all $s \in S$.
    \end{enumerate}
    The element $t$ is denoted by $t=\sup(S)$ and is called the \textbf{supremum} (least upper bound) of $S$.
  \end{define}    
    \begin{enumerate}[label=(\alph*)]
      \item For all \(x\in\setR\) there exists \(n\in\setZ\) with \(n>x\).
      \begin{subprove}
        Let \(x \in \setR\) be arbitrary. Suppose that 
        the set of integers \(\setZ\) is bounded above by some real number \(r \in \setR\)
        for the sake of contradiction. That is, 
        \[
          \exists r \in \setR \text{ such that } \forall n \in \setZ, n \le r.
        \]
        Then, by \nameref{def:completeness-axiom}, the nonempty set \(\setZ\) has 
        a least upper bound, denoted by \(s = \sup \setZ\).
        Consider \(s - 1\). Since \(s - 1 < s\), the number \(s - 1\) cannot 
        be an upper bound of \(\setZ\).
        Hence, by the definition of upper bound, 
        there exists an integer \(m \in \setZ\) such that
        \[
          m > s - 1.
        \]
        This implies that \(m + 1 > s\). By closure of the integers under addition, 
        it follows that \(m + 1 \in \setZ\) 
        which contradicts the fact that \(s\) is an upper bound of \(\setZ\).
        Therefore, by contradiction, it is true that 
        for any real number \(x \in \setR\), there exists an integer \(n \in \setZ\) with \(n > x\).
      \end{subprove}
      \item For all \(\varepsilon>0\) there exists \(n\in\setN\) with \(1/n<\varepsilon\).
      \begin{subprove}
        Let \(\varepsilon \in \setR\) be such that \(\varepsilon > 0\).
        It will be proved that there exists \(n \in \setN\) such that \(\frac{1}{n} < \varepsilon\).
        
        Since \(\varepsilon > 0\), the real number \(\frac{1}{\varepsilon}\) is well-defined and positive.
        By part (a), applied to the real number \(x = \frac{1}{\varepsilon}\), there exists an integer \(m \in \setZ\) such that
        \[
          m > \frac{1}{\varepsilon}.
        \]
        Since \(\frac{1}{\varepsilon} > 0\) and \(m > \frac{1}{\varepsilon}\), it follows that \(m > 0\).
        Therefore, \(m \in \setN\).
        Let \(n = m\). Then \(n \in \setN\) and \(n > \frac{1}{\varepsilon} > 0\).
        
        Since \(n > \frac{1}{\varepsilon} > 0\), the inequality may be inverted, giving
        \[
          0 < \frac{1}{n} < \varepsilon.
        \]
        Therefore, there exists \(n \in \setN\) such that \(\frac{1}{n} < \varepsilon\).
      \end{subprove}
      \item For all \(x\in\setR\) there exists \(n\in\setZ\) such that \(n\le x<n+1\).
      \begin{subprove}
        Let \(x \in \setR\) be arbitrary.
        It will be proved that there exists \(n \in \setZ\) such that \(n \le x < n + 1\).
        
        Consider the set
        \[
          A := \{ k \in \setZ : k \le x \}.
        \]
        The set \(A\) is nonempty because, for example, \(x - 1 \in \setR\), and by part (a) there exists \(m \in \setZ\) 
        such that \(m > x - 1\), so \(m - 1 \le x\) and hence \(m - 1 \in A\).
        Moreover, the set \(A\) is bounded from above by \(x\), since every \(k \in A\) satisfies \(k \le x\) by definition.
        
        By \nameref{def:completeness-axiom}, the nonempty set \(A\), which is bounded from above, has a least upper bound.
        Denote this least upper bound by \(s = \sup A\).
        It will be shown that \(s \in A\) and that \(s\) is an integer.
        
        Since \(A \subseteq \setZ\) and \(A\) is nonempty, all elements of \(A\) are integers.
        Because \(s = \sup A\) is the least upper bound of \(A\), for every \(\delta > 0\) there exists \(a \in A\) such that
        \[
          s - \delta < a \le s.
        \]
        In particular, taking \(\delta = 1\), there exists \(a \in A\) such that
        \[
          s - 1 < a \le s.
        \]
        Since \(a \in A\), it follows that \(a \in \setZ\).
        Now \(s - 1 < a \le s\) implies
        \[
          s - 1 < a \quad\text{and}\quad a \le s.
        \]
        Because \(a\) is an integer and \(a \le s\), it must be that \(a \le s < a + 1\).
        Thus \(s\) lies in the interval \([a, a+1)\), which implies \(a \le x\) (since \(a \in A\)) and \(s \ge a\).
        But \(s\) is the least upper bound of \(A\), and \(a \in A\) implies \(a \le s\).
        In particular, \(a\) is the greatest integer in \(A\).
        
        Set \(n := a\).
        Then \(n \in \setZ\) and \(n \le x\) holds because \(n \in A\) by definition.
        It remains to show that \(x < n + 1\).
        Since \(n\) is the greatest integer less than or equal to \(x\), the integer \(n + 1\) cannot belong to \(A\).
        Thus \(n + 1 \nleq x\), i.e., \(n + 1 > x\).
        Therefore,
        \[
          n \le x < n + 1.
        \]
        Hence, such an integer \(n\) exists for the given real number \(x\).
      \end{subprove}
      \item For all \(x\in\setR\) and \(N\in\setN\) there exists \(n\in\setZ\) such that \(\dfrac{n}{N}\le x<\dfrac{n+1}{N}\).
      \begin{subprove}
        Let \(x \in \setR\) and \(N \in \setN\) be given.
        It will be proved that there exists \(n \in \setZ\) such that
        \[
          \frac{n}{N} \le x < \frac{n+1}{N}.
        \]
        
        Consider the real number \(Nx\).
        By part (c), applied to \(Nx\), there exists an integer \(k \in \setZ\) such that
        \[
          k \le Nx < k + 1.
        \]
        Define \(n := k\).
        Then \(n \in \setZ\), and dividing the inequality by \(N > 0\) gives
        \[
          \frac{n}{N} \le x < \frac{n+1}{N}.
        \]
        Therefore, such an integer \(n\) exists for the given \(x\) and \(N\).
      \end{subprove}
      \item For all \(x\in\setR\) and \(\varepsilon>0\) there exists a rational \(q\in\setQ\) with \(|x-q|<\varepsilon\).
      \begin{subprove}
        Let \(x \in \setR\) and \(\varepsilon \in \setR\) with \(\varepsilon > 0\) be given.
        It will be proved that there exists a rational number \(q \in \setQ\) such that \(|x - q| < \varepsilon\).
        
        By part (b), there exists a positive integer \(N \in \setN\) such that
        \[
          \frac{1}{N} < \varepsilon.
        \]
        Now part (d) is applied to the real number \(x\) and this positive integer \(N\).
        By part (d), there exists an integer \(n \in \setZ\) such that
        \[
          \frac{n}{N} \le x < \frac{n+1}{N}.
        \]
        Define
        \[
          q := \frac{n}{N}.
        \]
        Then \(q\) is a rational number since \(n \in \setZ\) and \(N \in \setN\) with \(N > 0\).
        
        From the inequality above, it follows that
        \[
          0 \le x - q < \frac{1}{N}.
        \]
        Hence,
        \[
          |x - q| = x - q < \frac{1}{N} < \varepsilon.
        \]
        Therefore, there exists a rational number \(q \in \setQ\) such that \(|x - q| < \varepsilon\).
      \end{subprove}
    \end{enumerate}

  \item Prove that for all \(a,b\in\setR\)
  \[
    \max\{a,b\}=\frac{a+b+|a-b|}{2}
    \quad\text{and}\quad
    \min\{a,b\}=-\max\{-a,-b\}=\frac{a+b-|a-b|}{2}.
  \]
  \begin{subprove}
    Let \(a, b \in \setR\).
    It will be proved that \(\max\{a,b\} = \frac{a+b+|a-b|}{2}\) and 
    \(\min\{a,b\} = -\max\{-a,-b\} = \frac{a+b-|a-b|}{2}\).
    
    First, it will be shown that \(\max\{a,b\} = \frac{a+b+|a-b|}{2}\).
    By the definition of absolute value, two cases are considered.
    
    \textbf{Case 1:} Suppose \(a \ge b\).
    Then \(|a-b| = a-b\).
    Therefore,
    \[
      \frac{a+b+|a-b|}{2} = \frac{a+b+(a-b)}{2} = \frac{2a}{2} = a = \max\{a,b\}.
    \]
    
    \textbf{Case 2:} Suppose \(a < b\).
    Then \(|a-b| = -(a-b) = b-a\).
    Therefore,
    \[
      \frac{a+b+|a-b|}{2} = \frac{a+b+(b-a)}{2} = \frac{2b}{2} = b = \max\{a,b\}.
    \]
    
    Therefore, in both cases, \(\max\{a,b\} = \frac{a+b+|a-b|}{2}\).
    
    Next, it will be shown that \(\min\{a,b\} = -\max\{-a,-b\} = \frac{a+b-|a-b|}{2}\).
    
    First, it is established that \(\min\{a,b\} = -\max\{-a,-b\}\).
    By the definition of minimum, \(\min\{a,b\} = a\) if \(a \le b\), and \(\min\{a,b\} = b\) if \(b \le a\).
    By the definition of maximum, \(\max\{-a,-b\} = -a\) if \(-a \ge -b\), i.e., if \(a \le b\),
    and \(\max\{-a,-b\} = -b\) if \(-b \ge -a\), i.e., if \(b \le a\).
    Therefore, if \(a \le b\), then \(\min\{a,b\} = a = -(-a) = -\max\{-a,-b\}\).
    If \(b \le a\), then \(\min\{a,b\} = b = -(-b) = -\max\{-a,-b\}\).
    Therefore, \(\min\{a,b\} = -\max\{-a,-b\}\).
    
    Now it will be shown that \(-\max\{-a,-b\} = \frac{a+b-|a-b|}{2}\).
    By the previous result, \(\max\{-a,-b\} = \frac{(-a)+(-b)+|(-a)-(-b)|}{2} = \frac{-(a+b)+|b-a|}{2}\).
    Since \(|b-a| = |a-b|\), it follows that
    \[
      \max\{-a,-b\} = \frac{-(a+b)+|a-b|}{2}.
    \]
    Therefore,
    \[
      -\max\{-a,-b\} = -\frac{-(a+b)+|a-b|}{2} = \frac{a+b-|a-b|}{2}.
    \]
    
    Therefore, \(\min\{a,b\} = -\max\{-a,-b\} = \frac{a+b-|a-b|}{2}\).
    
    Therefore, for all \(a,b \in \setR\),
    \[
      \max\{a,b\} = \frac{a+b+|a-b|}{2} \quad\text{and}\quad
      \min\{a,b\} = -\max\{-a,-b\} = \frac{a+b-|a-b|}{2}.
    \]
  \end{subprove}

    
  \item Prove the Archimedean property:

  \begin{define}{Archimedean Property}[def:archimedean-property][red]
    For any real numbers \(a, b\) with \(a>0\) and \(b>0\), there exists 
    a natural number \(n\) such that 
    \[
    na > b.
    \]
    \textbf{Note}: This property states that the natural numbers are unbounded in the real numbers.
  \end{define}
  \begin{subprove}
    Suppose that there exist \(a, b \in \setR\) such that \(na \le b\) for all \(n \in \setN\) 
    for the sake of contradiction.
    Then, \(b\) is an upper bound for the set \(S = \set{na : n \in \setN}\).
    Hence, \(S\) has a supremum \(t=\sup(S) \in \setR\) by \nameref{def:completeness-axiom}.

    Since $a>0$, $t - a < t < t + a$. Hence, $t-a$ is not an upper bound for \(S\).
    Therefore, there exists \(n_0 \in \setN\) such that \(n_0 a > t - a\). But then, 
    \[
      t = t - a + a < n_0 a + a = (n_0 + 1)a,
    \]
    which contradicts the fact that \(t\) is an upper bound for \(S\).
    Therefore, by contradiction, the Archimedean Property is true.
  \end{subprove}

  \item Explain why \(\setQ\) is dense in \(\setR\):
    If \(a,b\in\setR\) and \(a<b\), then there exists \(r\in\setQ\) with \(a<r<b\).
    \begin{subprove}
      Let \(a, b \in \setR\) such that \(a < b\).
      It will be proved that there exists \(r \in \setQ\) such that \(a < r < b\).
      
      Since \(a < b\), it follows that \(b - a > 0\).
      By the \nameref{def:archimedean-property}, applied to \(b - a > 0\) and \(1 > 0\), 
      there exists \(n \in \setN\) such that \(n(b - a) > 1\).
      Therefore, \(b - a > \frac{1}{n}\), which implies \(\frac{1}{n} < b - a\).
      
      Consider the set of rational numbers of the form \(\frac{m}{n}\) where \(m \in \setZ\).
      Since the integers are unbounded in the real numbers (by the \nameref{def:archimedean-property}), 
      there exists an integer \(m\) such that \(m > na\).
      Let \(m_0\) be the smallest integer such that \(m_0 > na\).
      Then \(m_0 - 1 \le na < m_0\), which implies \(\frac{m_0 - 1}{n} \le a < \frac{m_0}{n}\).
      
      It will be shown that \(\frac{m_0}{n} < b\).
      Since \(m_0 > na\), it follows that \(\frac{m_0}{n} > a\).
      Suppose, for the sake of contradiction, that \(\frac{m_0}{n} \ge b\).
      Then \(\frac{m_0}{n} \ge b > a\), so \(\frac{m_0}{n} - a \ge b - a > \frac{1}{n}\).
      This implies \(\frac{m_0 - na}{n} > \frac{1}{n}\), so \(m_0 - na > 1\).
      Therefore, \(m_0 - 1 > na\), which contradicts that \(m_0\) is the smallest integer greater than \(na\).
      
      Therefore, \(a < \frac{m_0}{n} < b\).
      Since \(m_0 \in \setZ\) and \(n \in \setN\), it follows that \(r = \frac{m_0}{n} \in \setQ\).
      
      Therefore, there exists \(r \in \setQ\) such that \(a < r < b\).
    \end{subprove}
\end{enumerate}




























%%%%%%%%%%%%%%%%%%%%%%%%%%%%%%%%%%%%%%%%%%%%%%%%%%%%%%%%%%%%%%%
% Analysis 1 - HW
%%%%%%%%%%%%%%%%%%%%%%%%%%%%%%%%%%%%%%%%%%%%%%%%%%%%%%%%%%%%%%%
\newpage\handout[true]
{\textsc{Math 2106-B: Foundations of Mathematical Proof}}{Fall 2025}
{Homework 8: Completeness and Order in $\setR$}
{\textit{Prof.} \href{https://ceheitsch.github.io/webpage/}{Dr. Christine Heitsch}}{Due Date: November 20, 2025}
{Math 2106-B: HW8}
{\large
  \noindent
  \textbf{Name:} \HandoutAuthor \smallskip \\

  \noindent
  \textbf{GTID:} 903743506 \smallskip \\

  \noindent
  %%% List anyone with whom you discussed the problems
  \textbf{Collaborators:} None. This material reflects independent preparation. \smallskip \\

  \noindent
  %%% List all resources used OTHER than the textbook, lecture notes, 
  %%% webwork problems, ICA questions, and previous homework assignments.
  \textbf{Outside resources:} This work was informed by MATH 2106 course materials 
  (ICAs: in-class activities and WebWork contents) from \HandoutInstitution\ (the course textbook is 
  Chartrand \parencite{chartrand_mathematical_2018}).   
  Outside references specifically include
  textbooks written by 
  Grimaldi \parencite{grimaldi_discrete_2004} (used in MATH 3012) and 
  Rosen \parencite{rosen_discrete_2019} (used in CS 2050). \smallskip \\

  \noindent
  \textbf{Notice}: The answers contain cross-references throughout this document as clickable hyperlinks in most PDF viewers. However, the Gradescope PDF view might not support clickable hyperlinks.
}
\printbibliography[
  heading=bibintoc,   
  title={References}     % change “References” text if needed
]
% `heading=bibnumbered` if you prefer it numbered chapter (in \tableofcontents)
% `heading=bibintoc`    if you prefer it unnumbered but still in the ToC
\newpage
%%%%%%%%%%%%%%%%%%%%%%%%%%%%%%%%%%%%%%%%%%%%%%%%%%%%%%%%%%%%%%%%%%%%%%%
% Problems Section
%%%%%%%%%%%%%%%%%%%%%%%%%%%%%%%%%%%%%%%%%%%%%%%%%%%%%%%%%%%%%%%%%%%%%%%
\begin{enumerate}
  \item Let $a_1, a_2, \ldots, a_n \in \setR$ for $n \in \setZ$ with $n \geq 3$. 
  Prove that $\abs{\sum_{i=1}^n a_i} \leq \sum_{i=1}^n\abs{a_i}$.
  \begin{define}{Mathematical Induction}[def:mathematical-induction][red]
    \textbf{Mathematical induction} is a method of proof used to prove statements about natural numbers.
    To prove that a statement \(P(n)\) is true for all natural numbers \(n \geq n_{0}\), where \(n_{0}\) is some base value, it suffices to show that
    \begin{enumerate}
      \item \textbf{Base case:} \(P(n_{0})\) is true, and
      \item \textbf{Inductive step:} For any \(k \geq n_{0}\), if \(P(k)\) is true, then \(P(k+1)\) is true.
    \end{enumerate}
  \end{define}
  \begin{subprove}
    Let $P(n)$ be the proposition that 
    $\abs{\sum_{i=1}^n a_i} \leq \sum_{i=1}^n\abs{a_i}$ 
    for all $a_1, a_2, \ldots, a_n \in \setR$ with $n \geq 2$.
    \textbf{Base case:} For $n = 2$, the \nameref{def:triangle-inequality} states that for all $a_1, a_2 \in \setR$,
    \[
      |a_1 + a_2| \leq |a_1| + |a_2|.
    \]
    Therefore, $P(2)$ is true.
    
    \textbf{Inductive step:} Suppose that $P(k)$ is true for some $k \geq 2$ 
    for the sake of direct proof. 
    Let $a_1, a_2, \ldots, a_{k+1} \in \setR$, then
    \[
      \abs{\sum_{i=1}^{k+1} a_i} = \abs{\sum_{i=1}^k a_i + a_{k+1}} \leq \abs{\sum_{i=1}^k a_i} + |a_{k+1}|.
    \]
    by the \nameref{def:triangle-inequality} for two terms, then by the inductive hypothesis,
    \[
      \abs{\sum_{i=1}^k a_i} \leq \sum_{i=1}^k\abs{a_i}.
    \]
    Therefore,
    \[
      \abs{\sum_{i=1}^{k+1} a_i} \leq \sum_{i=1}^k\abs{a_i} + |a_{k+1}| = \sum_{i=1}^{k+1}\abs{a_i}.
    \]
    Thus, $P(k+1)$ is true.
    By \nameref{def:mathematical-induction}, it follows that $P(n)$ is true for all $n \geq 2$.
    Therefore, $\abs{\sum_{i=1}^n a_i} \leq \sum_{i=1}^n\abs{a_i}$ for all $a_1, a_2, \ldots, a_n \in \setR$ with $n \geq 3$.
  \end{subprove}





  \item Let $a, b \in \setR$. Prove that if $a \leq b_1$ for all real numbers $b_1>b$, then $a \leq b$.
  \begin{subprove}
    Suppose $a > b$ for the sake of contradiction with the given conditions.
    Then, $a + b > 2b$ by adding $b$ to both sides. Furthermore, $\dfrac{a+b}{2} > b$ by dividing both sides by 2.

    Let $b_1 = \dfrac{a+b}{2}$ by the previous result.
    Then, $a \leq \dfrac{a+b}{2}$ by the assumption that $a \leq b_1$ for all real numbers $b_1 > b$.
    It follows that $2a \leq a + b$ by multiplying both sides by 2.
    Furthermore, $a \leq b$ by subtracting $a$ from both sides, which contradicts the assumption that $a > b$.

    Therefore, by contradiction, if $a \leq b_1$ for all real numbers $b_1>b$, then $a \leq b$.
  \end{subprove}




  

  

  \newpage

  \item The \textit{completeness axiom} for the real numbers says that a nonempty subset 
  $S$ of $\setR$ that is bounded from above has a least upper bound; that is, 
  if $S \subset \setR$, $S \neq \emptyset$, and there exists $b \in \setR$ 
  such that $t \leq b$ for all $t \in S$, then $\sup S$ exists and is a real number. \\
  Let $S \subseteq \setR$ be nonempty and bounded from below, and let $-S=\{-s \mid s \in S\}$.
  \begin{enumerate}
    \item Prove that $\sup (-S)$ exists.
    \begin{subprove}
      Since $S$ is nonempty, there exists $s \in S$.
      Thus, $-s \in -S$, so $-S$ is \underbar{nonempty}.
      Since $S$ is bounded from below, there exists $l \in \setR$ such that $l \leq s$ for all $s \in S$.
      By multiplying both sides by $-1$, it follows that $-s \leq -l$ for all $-s \in -S$.
      Thus, $-l$ is an \underbar{upper bound} for $-S$.      
      Since $-S$ is nonempty and bounded above, by \nameref{def:completeness-axiom}, $\sup(-S)$ exists.
    \end{subprove}
    \item Let $s_0=\sup (-S)$. Prove that $-s_0 \leq s$ for all $s \in S$.
    \begin{subprove}
      Let $s_0 = \sup(-S)$ and let $s \in S$.      
      Since $s_0$ is the supremum of $-S$, by definition, $-s \leq s_0$ for all $-s \in -S$.
      By multiplying both sides by $-1$, it follows that $s \geq -s_0$.
      Thus, $-s_0 \leq s$.
      
      Therefore, $-s_0 \leq s$ for all $s \in S$.
    \end{subprove}
    \item Let $s_0=\sup (-S)$. Prove that if $t \leq s$ for all $s \in S$, then $t \leq -s_0$.
    \begin{subprove}
      Let $s_0 = \sup(-S)$ and suppose that $t \leq s$ for all $s \in S$ for the sake of direct proof.      
      Since $t \leq s$ for all $s \in S$, by multiplying both sides by $-1$, it follows that $-t \geq -s$ for all $-s \in -S$.
      This means that $-t$ is an \underbar{upper bound} for $-S$.
      Since $s_0 = \sup(-S)$ is the supremum of $-S$ by definition, it follows that $s_0 \leq -t$.
      By multiplying both sides by $-1$, it follows that $-s_0 \geq t$.
      Therefore, if $t \leq s$ for all $s \in S$, then $t \leq -s_0$.
    \end{subprove}


    \item Prove $\inf S=-\sup (-S)$.
    \begin{define}{Infimum (Greatest Lower Bound)}[def:infimum][red]
      Let $S\subseteq\setR$ be nonempty and bounded below.
      Then there exists an element $m\in\setR$ such that
      \begin{enumerate}
        \item $m\le s$ for all $s\in S$, and
        \item if $y\in\setR$ satisfies $y\le s$ for all $s\in S$ then $y\le m$.
      \end{enumerate}
      The element $m$ is denoted $m=\inf(S)$ and is called the \textbf{infimum} (greatest lower bound) of $S$.
        \begin{subprove}
          Let $S\subseteq\mathbb R$ be nonempty and bounded below. 
          Let $-S=\{-s:s\in S\}$, then $-S$ is nonempty and bounded above, 
          hence there exists $t=\sup(-S)$ by \nameref{def:completeness-axiom} and $m=-t$. \proofstep
          For any $s\in S$, $-s\le t$ and $t \le u$ for all upper bounds, $u$, by definition. 
          Thus, $m=-t\le s$ and $-u\le -t=m$ for all lower bounds, $-u$, of $S$. \proofstep      
          Therefore $\inf(S)=m=-\sup(-S)$.
        \end{subprove}
    \end{define}
    \begin{subprove}
      Let $s_0 = \sup(-S)$.
      By part (b), $-s_0 \leq s$ for all $s \in S$.
      By part (c), if $t \leq s$ for all $s \in S$, then $t \leq -s_0$.
      Thus, $-s_0$ satisfies both conditions of the definition of infimum of $S$ as follows.
      \begin{enumerate}
        \item $-s_0 \leq s$ for all $s \in S$, and
        \item if $t \leq s$ for all $s \in S$, then $t \leq -s_0$.
      \end{enumerate}
      Therefore, $\inf S = -s_0 = -\sup(-S)$.
    \end{subprove}
    \item \label{prob:3e-greatest-lower-bound} Conclude that every nonempty set of real numbers that is bounded from below has a greatest lower bound.
    \begin{subprove}
      Let $S \subseteq \setR$ be nonempty and bounded from below.
      It will be proved that $S$ has a \underbar{greatest lower bound}.
      By part (a), $\sup(-S)$ exists.
      By part (d), $\inf S = -\sup(-S)$.
      Thus, $\inf S$ exists and is a real number.
      
      Therefore, every nonempty set of real numbers that is bounded from below has a greatest lower bound.
    \end{subprove}
  \end{enumerate}










  \newpage
  
  \item Let $S$ be a nonempty subset of $\setR$ that is bounded from above. 
  Prove that if $\sup S$ belongs to $S$, then $\sup S=\max S$.
  \begin{subprove}
    Let $S$ be a nonempty subset of $\setR$ that is bounded from above, and suppose $\sup S \in S$.
    
    Since $\sup S$ is the supremum of $S$, by definition, $s \leq \sup S$ for all $s \in S$.
    Since $\sup S \in S$, it follows that $\sup S$ is an element of $S$ that is greater than or equal to all other elements of $S$.
    Thus, $\sup S$ is the \underbar{maximum} element of $S$.
    
    Therefore, if $\sup S$ belongs to $S$, then $\sup S = \max S$.
  \end{subprove}





  \item Let $S$ be a nonempty bounded 
  (that is, bounded from above and bounded from below) 
  subset of $\setR$. \textit{You may assume that Problem \ref{prob:3e-greatest-lower-bound} is true}.
  \begin{enumerate}
    \item Prove that $\inf S \leq \sup S$.
    \begin{subprove}
      Let $S$ be a nonempty bounded subset of $\setR$.
      
      Since $S$ is nonempty, there exists $s \in S$.
      By definition of infimum, $\inf S \leq s$ for all $s \in S$.
      By definition of supremum, $s \leq \sup S$ for all $s \in S$.
      Thus, for any $s \in S$, it follows that $\inf S \leq s \leq \sup S$.
      Therefore, $\inf S \leq \sup S$.
    \end{subprove}
    \item What can be said about $S$ if $\inf S=\sup S$? Justify your answer.
    \begin{subanswer}
      If $\inf S = \sup S$, then $S$ is a singleton set, that is, $S$ contains exactly one element.
      
      Let $c = \inf S = \sup S$.
      By definition of infimum, $c \leq s$ for all $s \in S$.
      By definition of supremum, $s \leq c$ for all $s \in S$.
      Thus, $s = c$ for all $s \in S$.
      Since $S$ is nonempty, there exists at least one element in $S$, and since all elements equal $c$, it follows that $S = \{c\}$.
    \end{subanswer}
  \end{enumerate}






  \newpage

  \item Let $S$ and $T$ be nonempty bounded subsets of $\setR$.
  \begin{enumerate}
    \item Prove that if $S \subseteq T$, then $\inf T \leq \inf S \leq \sup S \leq \sup T$.
    \begin{subprove}
      Let $S$ and $T$ be nonempty bounded subsets of $\setR$ such that $S \subseteq T$.
      \begin{enumerate}
        \item 
        By the definition of subset, for all $s \in S$, $s \in T$.
        By the definition of infimum, $\inf T \leq t$ for all $t \in T$.
        It follows that $\inf T \leq s$ for all $s \in S$ since $s \in T$.
        Thus, $\inf T$ is a lower bound of $S$.
  
        By the definition of infimum, since $\inf S$ is the greatest lower bound of $S$, it follows that $\inf T \leq \inf S$.
        \item
        By the definition of subset, for all $s \in S$, $s \in T$.
        By the definition of supremum, $t \leq \sup T$ for all $t \in T$.
        It follows that $s \leq \sup T$ for all $s \in S$ since $s \in T$.
        Thus, $\sup T$ is an upper bound of $S$.
  
        By the definition of supremum, since $\sup S$ is the least upper bound of $S$, it follows that $\sup S \leq \sup T$.
        \item
        By Problem 5(a), $\inf S \leq \sup S$.
      \end{enumerate}
      Therefore, if $S \subseteq T$, then $\inf T \leq \inf S \leq \sup S \leq \sup T$.
    \end{subprove}
    \item Prove that $\sup (S \cup T)=\max \{\sup S, \sup T\}$ without assuming $S \subseteq T$.
    \begin{subprove}
      Let $S$ and $T$ be nonempty bounded subsets of $\setR$, and let $M = \max \{\sup S, \sup T\}$.
      
      By the definition of union, $S \cup T = \sBUILD{x \in \setR}{x \in S \text{ or } x \in T}$.
      For any $x \in S \cup T$, by the definition of union, either $x \in S$ or $x \in T$.
      If $x \in S$, then by the definition of supremum, $x \leq \sup S$.
      If $x \in T$, then by the definition of supremum, $x \leq \sup T$.
      Since $M = \max \{\sup S, \sup T\}$, it follows that $\sup S \leq M$ and $\sup T \leq M$.
      Therefore, for any $x \in S \cup T$, $x \leq M$.
      By the definition of upper bound, $M$ is an \underbar{upper bound} for $S \cup T$.
      
      Suppose that there exists an upper bound $u$ of $S \cup T$ such that $u < M$ for the sake of contradiction.
      By the definition of maximum, $M = \max \{\sup S, \sup T\}$ implies that either $M = \sup S$ or $M = \sup T$.
      Without loss of generality, assume $M = \sup S$.
      Then, $u < \sup S$.
      Since $u$ is an upper bound of $S \cup T$, by the definition of upper bound, $x \leq u$ for all $x \in S \cup T$.
      Since $S \subseteq S \cup T$, it follows that $x \leq u$ for all $x \in S$.
      Thus, $u$ is an upper bound of $S$.
      However, by the definition of supremum, since $\sup S$ is the least upper bound of $S$ and $u < \sup S$, this contradicts the fact that $\sup S$ is the least upper bound.
      Therefore, no such $u$ exists, and $M$ is the \underbar{least upper bound} of $S \cup T$.
      
      Therefore, $\sup (S \cup T) = M = \max \{\sup S, \sup T\}$.
    \end{subprove}
  \end{enumerate}












  \newpage
  \item \textbf{Complete but do not hand in.} Prove the following properties of the absolute value.
  \begin{enumerate}
    \item For all $a, b \in \setR$, $||a|-|b|| \leq |a-b|$.
    \begin{subprove}
      Let $a, b \in \setR$.
      
      By the \nameref{def:triangle-inequality}, $|a| = |(a-b) + b| \leq |a-b| + |b|$.
      It follows that $|a| - |b| \leq |a-b|$.
      By the \nameref{def:triangle-inequality}, $|b| = |(b-a) + a| \leq |b-a| + |a| = |a-b| + |a|$.
      It follows that $|b| - |a| \leq |a-b|$.
      Therefore, $-(|a| - |b|) = |b| - |a| \leq |a-b|$.
      Thus, $-|a-b| \leq |a| - |b| \leq |a-b|$.
      By the definition of absolute value, it follows that $||a| - |b|| \leq |a-b|$.
      
      Therefore, for all $a, b \in \setR$, $||a|-|b|| \leq |a-b|$.
    \end{subprove}
    \item If $a, x, \epsilon \in \setR$, then $|x-a|<\epsilon$ if and only if $a-\epsilon<x<a+\epsilon$.
    \begin{subprove}
      Let $a, x, \epsilon \in \setR$.
      It will be proved that $|x-a| < \epsilon$ if and only if $a-\epsilon < x < a+\epsilon$.
      
      Suppose that $|x-a| < \epsilon$.
      By the definition of absolute value, $|x-a| < \epsilon$ means that either $x-a < \epsilon$ and $x-a \geq 0$, or $-(x-a) < \epsilon$ and $x-a < 0$.
      If $x-a \geq 0$, then $x-a < \epsilon$, so $x < a+\epsilon$.
      Since $x-a \geq 0$, it follows that $x \geq a > a-\epsilon$.
      If $x-a < 0$, then $-(x-a) = a-x < \epsilon$, so $a-\epsilon < x$.
      Since $x-a < 0$, it follows that $x < a < a+\epsilon$.
      Therefore, in both cases, $a-\epsilon < x < a+\epsilon$.
      
      Conversely, suppose that $a-\epsilon < x < a+\epsilon$.
      It follows that $-\epsilon < x-a < \epsilon$.
      By the definition of absolute value, since $-\epsilon < x-a < \epsilon$, it follows that $|x-a| < \epsilon$.
      
      Therefore, $|x-a| < \epsilon$ if and only if $a-\epsilon < x < a+\epsilon$.
    \end{subprove}
    \item If $a, b, x, y \in \setR$ and $a<x<b$ with $a<y<b$, then $|y-x|<b-a$.
    \begin{subprove}
      Let $a, b, x, y \in \setR$ such that $a < x < b$ and $a < y < b$.
      
      By the definition of absolute value, $|y-x| = y-x$ if $y \geq x$, and $|y-x| = x-y$ if $y < x$.
      If $y \geq x$, then since $y < b$ and $x > a$, it follows that $|y-x| = y-x < b-a$.
      If $y < x$, then since $x < b$ and $y > a$, it follows that $|y-x| = x-y < b-a$.
      Therefore, in both cases, $|y-x| < b-a$.
      
      Therefore, if $a < x < b$ and $a < y < b$, then $|y-x| < b-a$.
    \end{subprove}
  \end{enumerate}




  \newpage

  \item \textbf{Complete but do not hand in.} Using the completeness axiom and the ordering of real numbers, prove each of the 
  following statements.
  \begin{enumerate}
    \item For all $x \in \setR$, there exists $n \in \setZ$ such that $n>x$.
    \begin{subprove}
      Let $x \in \setR$ be arbitrary.
    
      For the sake of contradiction, suppose that the set of integers $\setZ$ is bounded from above.
      Then, by the completeness axiom, the nonempty set $\setZ$ has a least upper bound; denote it by $s = \sup \setZ$.
    
      Consider the real number $s - 1$. Since $s - 1 < s$, the number $s - 1$ cannot be an upper bound of $\setZ$.
      Hence, by the definition of upper bound, there exists an integer $m \in \setZ$ such that
      \[
        m > s - 1.
      \]
      Since $s$ is an upper bound of $\setZ$, it follows that $m \leq s$.
      Therefore,
      \[
        s - 1 < m \leq s.
      \]
      Now consider $m + 1$. Since $m \in \setZ$, it follows that $m + 1 \in \setZ$.
      Moreover,
      \[
        m + 1 > s.
      \]
      This contradicts the fact that $s$ is an upper bound of $\setZ$, because an element of $\setZ$ has been found that is greater than $s$.
    
      Therefore, the assumption that $\setZ$ is bounded from above must be false.
      It follows that $\setZ$ is not bounded from above.
      By the definition of being unbounded above, for any real number $x \in \setR$, there exists an integer $n \in \setZ$ such that $n > x$.
    \end{subprove}
    
    \item For all $\epsilon \in \setR$ with $\epsilon>0$, there exists $n \in \setZ$ with $n>0$ such that $1 / n<\epsilon$.
    \begin{subprove}
      Let $\epsilon \in \setR$ be such that $\epsilon > 0$.
      It will be shown that there exists $n \in \setZ$ with $n > 0$ such that $\dfrac{1}{n} < \epsilon$.
    
      Since $\epsilon > 0$, the real number $\dfrac{1}{\epsilon}$ is well-defined and positive.
      By part (a), applied to the real number $x = \dfrac{1}{\epsilon}$, there exists an integer $n \in \setZ$ such that
      \[
        n > \frac{1}{\epsilon}.
      \]
      In particular, $n > 0$ holds as $\dfrac{1}{\epsilon} > 0$.
    
      Because $n > \dfrac{1}{\epsilon}$ and $n > 0$, the inequality may be inverted, giving
      \[
        0 < \frac{1}{n} < \epsilon.
      \]
      Thus, a positive integer $n$ has been found such that $\dfrac{1}{n} < \epsilon$.
    \end{subprove}
    
    \item For all $x \in \setR$, there exists $n \in \setZ$ such that $n \leq x<n+1$.
    \begin{subprove}
      Let $x \in \setR$ be arbitrary.
      It will be shown that there exists $n \in \setZ$ such that $n \leq x < n + 1$.
    
      Consider the set
      \[
        A := \{ k \in \setZ : k \leq x \}.
      \]
      The set $A$ is nonempty because, for example, $x - 1 \in \setR$, and by part (a) there exists $m \in \setZ$ such that $m > x - 1$, so $m - 1 \leq x$ and hence $m - 1 \in A$.
      Moreover, the set $A$ is bounded from above by $x$, since every $k \in A$ satisfies $k \leq x$ by definition.
    
      By the completeness axiom, the nonempty set $A$, which is bounded from above, has a least upper bound.
      Denote this least upper bound by $s = \sup A$.
      It will be shown that $s \in A$ and that $s$ is an integer.
    
      Since $A \subseteq \setZ$ and $A$ is nonempty, all elements of $A$ are integers.
      Because $s = \sup A$ is the least upper bound of $A$, for every $\delta > 0$ there exists $a \in A$ such that
      \[
        s - \delta < a \leq s.
      \]
      In particular, taking $\delta = 1$, there exists $a \in A$ such that
      \[
        s - 1 < a \leq s.
      \]
      Since $a \in A$, it follows that $a \in \setZ$.
      Now $s - 1 < a \leq s$ implies
      \[
        s - 1 < a \quad \text{and} \quad a \leq s.
      \]
      Because $a$ is an integer and $a \leq s$, it must be that $a \leq s < a + 1$.
      Thus $s$ lies in the interval $[a, a+1)$, which implies $a \leq x$ (since $a \in A$) and $s \geq a$.
      But $s$ is the least upper bound of $A$, and $a \in A$ implies $a \leq s$.
      In particular, $a$ is the greatest integer in $A$.
    
      Set $n := a$.
      Then $n \in \setZ$ and $n \leq x$ holds because $n \in A$ by definition.
      It remains to show that $x < n + 1$.
      Since $n$ is the greatest integer less than or equal to $x$, the integer $n + 1$ cannot belong to $A$.
      Thus $n + 1 \nleq x$, i.e., $n + 1 > x$.
      Therefore,
      \[
        n \leq x < n + 1.
      \]
      Hence, such an integer $n$ exists for the given real number $x$.
    \end{subprove}
    
    \item For all $x \in \setR$ and $N \in \setZ$ with $N>0$, there exists $n \in \setZ$ such that $n / N \leq x<(n+1) / N$.
    \begin{subprove}
      Let $x \in \setR$ and $N \in \setZ$ with $N > 0$ be given.
      It will be shown that there exists $n \in \setZ$ such that
      \[
        \frac{n}{N} \leq x < \frac{n+1}{N}.
      \]
    
      Consider the real number $Nx$.
      By part (c), applied to $Nx$, there exists an integer $k \in \setZ$ such that
      \[
        k \leq Nx < k + 1.
      \]
      Define $n := k$.
      Then $n \in \setZ$, and dividing the inequality by $N > 0$ gives
      \[
        \frac{n}{N} \leq x < \frac{n+1}{N}.
      \]
      Thus, such an integer $n$ exists for the given $x$ and $N$.
    \end{subprove}
    
    \item For all $x, \epsilon \in \setR$ with $\epsilon>0$, there exists a rational number $q \in \setQ$ such that $|x-q|<\epsilon$.
    \begin{subprove}
      Let $x \in \setR$ and $\epsilon \in \setR$ with $\epsilon > 0$ be given.
      It will be shown that there exists a rational number $q \in \setQ$ such that $|x - q| < \epsilon$.
    
      By part (b), there exists a positive integer $N \in \setZ$ such that
      \[
        \frac{1}{N} < \epsilon.
      \]
      Now part (d) is applied to the real number $x$ and this positive integer $N$.
      By part (d), there exists an integer $n \in \setZ$ such that
      \[
        \frac{n}{N} \leq x < \frac{n+1}{N}.
      \]
      Define
      \[
        q := \frac{n}{N}.
      \]
      Then $q$ is a rational number since $n \in \setZ$ and $N \in \setZ$ with $N > 0$.
    
      From the inequality above, it follows that
      \[
        0 \leq x - q < \frac{1}{N}.
      \]
      Hence,
      \[
        |x - q| = x - q < \frac{1}{N} < \epsilon.
      \]
      Therefore, a rational number $q$ has been found such that $|x - q| < \epsilon$.
    \end{subprove}
    
  \end{enumerate}


\end{enumerate}






































%%%%%%%%%%%%%%%%%%%%%%%%%%%%%%%%%%%%%%%%%%%%%%%%%%%%%%%%%%%%%%%
% Analysis 2
%%%%%%%%%%%%%%%%%%%%%%%%%%%%%%%%%%%%%%%%%%%%%%%%%%%%%%%%%%%%%%%




\newpage





\begin{enumerate}
  \item Let \((P_n)\) be the $n$-th integer in the decimal expansion of \(\pi\), 
  what are the first five terms of the sequence \((P_n)_{n=1}^\infty\)?
  \begin{define}{Sequence}[def:sequence][red]
    A sequence of real numbers is a real-valued function
    defined on a domain \(D_m=\set{n\in\setZ | n \geq m}\) for some \(m \in \setZ\).

    The outputs are called the \textbf{terms} of the sequence and written as \(s_n\) for \(n \in D_m\).

    Sequences are written as \((s_n)_{n=m}^\infty\), \((s_n)_{n\in\setN}\), or just \((s_n)\)
    if the domain is arbitrary.
  \end{define}
  \begin{subanswer}
    The first five terms of the sequence \((P_n)_{n=1}^\infty\) are:
    \[
      P_1 = 1, \quad P_2 = 4, \quad P_3 = 1, \quad P_4 = 5, \quad P_5 = 9.
    \]
  \end{subanswer}
  \item Prove the following for \(n \in \setN\) from the definitions.
  \begin{define}{Convergence of a Sequence}[def:convergence-of-a-sequence][red]
    Let \((s_n)_{n\in\setN}\) be a sequence of \(\setR\).
    A sequence \((s_n)_{n\in\setN}\) in \(\setR\) \textbf{converges} to a limit \(s \in \setR\)
    if for every \(\varepsilon>0\) there exists \(N\in\setN\) such that 
    \[
      \abs{s_n - s} < \varepsilon \quad\text{for all } n > N.
    \]
    If \((s_n)_{n\in\setN}\) converges to \(s\), then
    \[
      \lim_{n\to\infty} s_n = s
    \]
    and $s$ is called the \textbf{limit} of the sequence.
  \end{define}
  \begin{enumerate}
    \item Prove that 
    \[
      \lim_{n\to\infty} \frac{1}{n} = 0.
    \]
    \begin{subprove}
      Let \(\varepsilon>0\) be given and 
      \(N=\ceil{\dfrac{1}{\varepsilon}}\).
      Then, for all \(n > N \ge \dfrac{1}{\varepsilon} > 0\), 
      it follows that \(\dfrac{1}{n} < \varepsilon\).
      Therefore,
      \[
        \abs{\frac{1}{n} - 0} = \frac{1}{n} < \varepsilon.
      \]
      Thus, by \nameref{def:convergence-of-a-sequence}, the sequence converges to \(0\).
    \end{subprove}





    \newpage

    \item \(\lim_{n\to\infty}\dfrac{(-1)^n}{n}=0\).
    \begin{subprove}
      Let \(\varepsilon>0\) be given and 
      $N=\ceil{\dfrac{1}{\varepsilon}}$.
      Then, for all \(n > N \ge \dfrac{1}{\varepsilon} > 0\), 
      it follows that \(\dfrac{1}{n} < \varepsilon\).
      Therefore,
      \[
        \abs{\dfrac{(-1)^n}{n}-0}=\dfrac{1}{n}<\varepsilon.
      \]
      Thus, by \nameref{def:convergence-of-a-sequence}, the sequence converges to \(0\).
    \end{subprove}

    \item \(\lim_{n\to\infty}\dfrac{1}{n^2}=0\).
    \begin{subprove}
      Let \(\varepsilon>0\) be given and 
      \(N=\ceil{\sqrt{\dfrac{1}{\varepsilon}}}\). 
      Then, for all \(n > N \ge \sqrt{\dfrac{1}{\varepsilon}} > 0\), 
      it follows that \(\dfrac{1}{n^2} < \varepsilon\).
      Therefore,
      \[
        \abs{\dfrac{1}{n^2}-0}=\dfrac{1}{n^2}<\varepsilon.
      \]
      Thus, by \nameref{def:convergence-of-a-sequence}, the sequence converges to \(0\).
    \end{subprove}


    \item Prove that \(\lim_{n\to\infty}\dfrac{4n^3+3n}{n^3-6}=4\).
    \begin{subprove}
      Let \(\varepsilon>0\) be given and 
      \(N=\max\set{3, \ceil{\sqrt{\dfrac{54}{\varepsilon}}}}\).
      Then, for all \(n > N \ge \sqrt{\dfrac{54}{\varepsilon}} > 0\), 
      it follows that \(\dfrac{54}{n^2} < \varepsilon\) and 
      \(n^3-6 > n^3/2 > 0\) since \(n > 3\).
      Therefore,
      \[
        \abs{\dfrac{4n^3+3n}{n^3-6}-4}
        =\abs{\dfrac{3n+24}{n^3-6}}
        <\dfrac{27n}{n^3/2}=\dfrac{54}{n^2}
        <\varepsilon.
      \]
      Thus, by \nameref{def:convergence-of-a-sequence}, the sequence converges to \(4\).
    \end{subprove}


    \item \(\lim_{n\to\infty}\dfrac{3n+1}{7n-4}=\dfrac{3}{7}\).
    \begin{subprove}
      Let \(\varepsilon>0\) be given and 
      \(N=\max\set{1, \ceil{\dfrac{1}{\varepsilon}}}\).
      Then, for all \(n > N \ge \dfrac{1}{\varepsilon}\), 
      it follows that \(\dfrac{1}{n} < \varepsilon\) and
      \(7n-4 > 3n > 0\) since \(n > 1\).
      Therefore,
      \[
        \abs{\dfrac{3n+1}{7n-4} - \dfrac{3}{7}}
        =\abs{\dfrac{7(3n+1) - 3(7n-4)}{7(7n-4)}}
        =\abs{\dfrac{19}{7(7n-4)}}
        <\dfrac{21}{7(3n)}
        <\varepsilon.
      \]
      Thus, by \nameref{def:convergence-of-a-sequence}, the sequence converges to \(\dfrac{3}{7}\).
    \end{subprove}


    \item Prove that \(\lim_{n\to\infty}\dfrac{n+6}{n^2-6}=0\) from the definition of sequence convergence.
    \begin{subprove}
      Let \(\varepsilon>0\) be given and \(N=\max\set{3, \ceil{\dfrac{14}{\varepsilon}}}\).
      Then, for all \(n > N \ge \dfrac{14}{\varepsilon} > 0\),
      it follows that \(\dfrac{14}{n} < \varepsilon\) and \(n^2-6 > n^2/2 > 0\) since \(n > 3\).
      Therefore,
      \[
        \abs{\dfrac{n+6}{n^2-6} - 0}
        <\dfrac{7n}{n^2 / 2} = \dfrac{14}{n}
        <\varepsilon.
      \]
      Thus, by \nameref{def:convergence-of-a-sequence}, the sequence converges to \(0\).
    \end{subprove}

    \newpage

    \item Let $p_n$ be the $n$-th digit in the decimal expansion of $\frac{23}{27}=0.851851\ldots$, 
    prove that $p_n$ does not converge.
    \begin{define}{Divergence of a Sequence}[def:divergence-of-a-sequence][red]
      A sequence \((s_n)_{n\in\setN}\) that does not converge is said to \textbf{diverge}
    \end{define}
    \begin{subprove}
      Suppose that \((p_n)\) converges to a limit \(p \in \setR\) for the sake of contradiction.
      Let \(\varepsilon = 0.1\). Then,
      by \nameref{def:convergence-of-a-sequence}, there exists \(N \in \setN\) such that 
      \[
        \abs{p_n - p} < 0.1 \quad\text{for all } n > N.
      \]
      However, the sequence \((p_n)\) alternates between the digits \(8\), \(5\), and \(1\).
      Thus, there must exist an \(n > N\) such that \(p_{n} = 8\) and \(p_{n+1} = 5\). Then,
      \[
        \abs{p_{n} - p} = \abs{8 - p} \quad\text{and}\quad \abs{p_{n+1} - p} = \abs{5 - p}.
      \]
      By the \nameref{def:triangle-inequality},
      \[
        3 = \abs{(8-p) - (5-p)} \leq \abs{8 - p} + \abs{5 - p} < 2 \cdot \varepsilon = 0.2
      \]
      which is a contradiction. Therefore, by contradiction,
      the sequence \((p_n)\) does not converge, and hence diverges.
    \end{subprove}

    \item The sequence \(a_n=(-1)^n\) does not converge.
    \begin{subprove}
      Suppose that \((a_n)\) converges to a limit \(a \in \setR\) for the sake of contradiction.
      Let \(\varepsilon = 0.1\). Then,
      by \nameref{def:convergence-of-a-sequence}, there exists \(N \in \setN\) such that
      \[
        \abs{a_n - a} < 0.1 \quad\text{for all } n > N.
      \]
      However, the sequence \((a_n)\) alternates between \(1\) and \(-1\).
      Thus, there must exist an \(n > N\) such that \(a_{n} = 1\) and \(a_{n+1} = -1\). Then,
      \[
        \abs{a_n - a} = \abs{1 - a} \quad\text{and}\quad \abs{a_{n+1} - a} = \abs{-1 - a}.
      \]
      By the \nameref{def:triangle-inequality},
      \[
        2 = \abs{(1-a) - (-1-a)} \leq \abs{1 - a} + \abs{-1 - a} < 2 \cdot \varepsilon = 0.2
      \]
      which is a contradiction. Therefore, by contradiction, the sequence \((a_n)\) does not converge, and hence diverges.
    \end{subprove}



    \newpage
      
    \item Give an example of a bounded divergent sequence other than \(((-1)^n)\).
    \begin{subanswer}
      The sequence \(s_n = \frac{n}{n+1}\cdot (-1)^n\) is bounded and divergent.
      To see that it is bounded, let \(M = 1\). Then, for all \(n \in \setN\),
      \[
        \abs{s_n} = \abs{\frac{n}{n+1}} \leq 1.
      \]
      Therefore, by \nameref{def:bounded-sequence}, the sequence \((s_n)\) is bounded.
      
      To see that it is divergent, suppose that \((s_n)\) converges to a limit \(a \in \setR\) for the sake of contradiction.
      Let \(\varepsilon = 0.1\). Then,
      by \nameref{def:convergence-of-a-sequence}, there exists \(N \in \setN\) such that
      \[
        \abs{s_n - a} < 0.1 \quad\text{for all } n > N.
      \]
      Let \(n > N\) be a large enough even number (at least \(2\)) such that 
      \(s_{n} = \frac{n}{n+1}\) and \(s_{n+1} = -\frac{n+1}{n+2}\).
      Then, by the \nameref{def:triangle-inequality},
      \[
        \abs{2 - \frac{1}{3} - \frac{1}{4}} < \abs{2 - \frac{1}{n+1} - \frac{1}{n+2}} = \abs{\frac{n}{n+1} + \frac{n+1}{n+2}}=\abs{s_n - s_{n+1}} \leq \abs{s_n - a} + \abs{s_{n+1} - a} < 2 \cdot \varepsilon = 0.2.
      \]
      which is a contradiction. Therefore, by contradiction, the sequence \((s_n)\) does not converge, and hence diverges.
    \end{subanswer}

    \item \(\lim_{n\to\infty}\sqrt{n}=+\infty\).
    \begin{define}{Divergence to Infinity}[def:divergence-to-infinity][red]
      Let \((s_n)_{n\in\setN}\) be a sequence in \(\setR\).
      \begin{itemize}
        \item The sequence \((s_n)\) is said to \textbf{diverge to infinity} (or \textbf{diverge to \(+\infty\)}) if,
          for every \(M\in\setR\), there exists \(N\in\setN\) such that
          \[
          s_n > M \quad\text{for all } n > N.
          \]
          In this case it is written \(\lim_{n\to\infty} s_n = +\infty\).
        \item The sequence \((s_n)\) is said to \textbf{diverge to negative infinity} (or \textbf{diverge to \(-\infty\)}) if,
          for every \(M\in\setR\), there exists \(N\in\setN\) such that
          \[
          s_n < M \quad\text{for all } n > N.
          \]
          In this case it is written \(\lim_{n\to\infty} s_n = -\infty\).
      \end{itemize}
      Such a sequence is not convergent to any finite real number.
    \end{define}

    \begin{subprove}
      Let \(M \in \setR\) be given and \(N=\ceil{M^{2}}\).
      Then, for all \(n > N \ge M^2 > 0\), it follows that \(\sqrt{n} > M\).
      Therefore,
      \[
        \sqrt{n} > M.
      \]
      Thus, by \nameref{def:divergence-to-infinity}, the sequence diverges to \(+\infty\).
    \end{subprove}

  \end{enumerate}

  \newpage

  \item Determine the limits of the following sequences \((s_n)_{n\in\setN}\), if they exist. Prove your answer.
  \begin{define}
    {Existence of Limit}[def:existence-of-limit][red]
    A sequence \((s_n)_{n\in\setN}\) in \(\setR\) \textbf{has a limit} if it either converges to a limit in \(\setR\)
    or diverges to \(+\infty\) or \(-\infty\).
  \end{define}
  \begin{enumerate}
    \item \(s_n=\dfrac{n}{n^2+1}\).
    \begin{subprove}
      Let \(\varepsilon>0\) be given and 
      \(N=\ceil{\dfrac{1}{\varepsilon}}\).
      Then, for all \(n > N \ge \dfrac{1}{\varepsilon} > 0\), 
      it follows that \(\dfrac{1}{n} < \varepsilon\).
      Since \(n^2+1 > n^2 > 0\), it follows that \(\dfrac{n}{n^2+1} < \dfrac{n}{n^2} = \dfrac{1}{n}\).
      Therefore,
      \[
        \abs{\dfrac{n}{n^2+1} - 0} = \dfrac{n}{n^2+1} < \dfrac{1}{n} < \varepsilon.
      \]
      Thus, by \nameref{def:convergence-of-a-sequence}, the sequence converges to \(0\).
      Therefore, by the \nameref{def:existence-of-limit}, the sequence \((s_n)\) has a limit.
    \end{subprove}
    \item \(s_n=\cos(n\pi/3)\).
    \begin{subprove}
      Suppose that \((s_n)\) converges to a limit \(s \in \setR\) for the sake of contradiction.
      Let \(\varepsilon = 0.1\). Then,
      by \nameref{def:convergence-of-a-sequence}, there exists \(N \in \setN\) such that
      \[
        \abs{s_n - s} < 0.1 \quad\text{for all } n > N.
      \]
      Let \(n\) be a multiple of 6 where \(n > N\) such that
      \[
        s_n = \cos(n\pi/3) = \cos(2k\pi) = 1 \quad\text{for some } k \in \setN
      \]
      and
      \[
        s_{n+3} = \cos((n+3)\pi/3) = \cos(2k\pi + \pi) = \cos(\pi) = -1 \quad\text{for some } k \in \setN.
      \]
      Then,
      \[
        \abs{s_n - s} = \abs{1 - s} \quad\text{and}\quad \abs{s_{n+3} - s} = \abs{-1-s} = \abs{s+1}.
      \]
      By the \nameref{def:triangle-inequality}, it follows that
      \[
        2 = \abs{(1-s) + (s+1)} \leq \abs{1 - s} + \abs{s+1} < 2 \cdot \varepsilon = 0.2.
      \]
      which is a contradiction. Therefore, by contradiction, the sequence \((s_n)\) does not converge. 
      Furthermore, it follows that \((s_n)\) is bounded since \(\abs{s_n} \leq 1\) 
      for all \(n \in \setN\) by definition of cosine.
      Therefore, by the \nameref{def:existence-of-limit}, the sequence \((s_n)\) does not have a limit.
    \end{subprove}
    \newpage

    \item \(s_n=(1/n)\sin n\).
    \begin{subprove}
      Let \(\varepsilon>0\) be given and 
      \(N=\ceil{\dfrac{1}{\varepsilon}}\).
      Then, for all \(n > N \ge \dfrac{1}{\varepsilon} > 0\), 
      it follows that \(\dfrac{1}{n} < \varepsilon\).
      Since \(|\sin n| \leq 1\) for all \(n \in \setN\) by definition of sine,
      it follows that \(\dfrac{1}{n}|\sin n| \leq \dfrac{1}{n}\).
      Therefore,
      \[
        \abs{\dfrac{1}{n}\sin n - 0} = \dfrac{1}{n}|\sin n| \leq \dfrac{1}{n} < \varepsilon.
      \]
      Thus, by \nameref{def:convergence-of-a-sequence}, the sequence converges to \(0\).
      Therefore, by the \nameref{def:existence-of-limit}, the sequence \((s_n)\) has a limit.
    \end{subprove}

    \item \(s_n=(-1)^n n\).
    \begin{subprove}
      Suppose that \((s_n)\) converges to \(s \in \setR\).
      Let \(\varepsilon = 0.1\). Then,
      by \nameref{def:convergence-of-a-sequence}, there exists \(N \in \setN\) such that
      \[
        \abs{s_n - s} < 0.1 \quad\text{for all } n > N.
      \]
      Let \(n > N\) be a large enough even number. Then,
      \[
        \abs{s_n - s_{n+1}} = \abs{(-1)^n n - (-1)^{n+1}(n+1)} = \abs{(-1)^n (n + (n+1))} = \abs{2n+1} = 2n+1.
      \]
      Then, by the \nameref{def:triangle-inequality},
      \[
        3 \le 2n+1 = \abs{s_n - s_{n+1}} \leq \abs{s_n - s} + \abs{s_{n+1} - s} < 0.1 + 0.1 = 0.2.
      \]
      which is a contradiction.
      Thus, \((s_n)\) does not converge to a limit in \(\setR\).
      Furthermore, suppose that \((s_n)\) diverges to \(+\infty\) for the sake of contradiction without loss of generality.
      Then, for \(M = 1\), there exists \(N \in \setN\) such that
      \[
        s_n > 1 \quad\text{for all } n > N.
      \]
      Let \(n > N\) be any possible odd number. Then,
      \[
        1 < s_n = (-1)^n n = -n < -1.
      \]
      which is a contradiction. Therefore, by contradiction, the sequence \((s_n)\) does not diverge 
      to \(+\infty\) or \(-\infty\).
      Therefore, by the \nameref{def:existence-of-limit}, the sequence \((s_n)\) does not have a limit.
    \end{subprove}
  \end{enumerate}


  \newpage

  \item Prove the following statement, Uniqueness of Limits.
  \begin{define}{Uniqueness of Limits}[lem:uniqueness-of-limits][red]
    A sequence in \(\setR\) has at most one limit.
  \end{define}
  \begin{subprove}
    Suppose a sequence \((s_n)_{n\in\setN}\) has two limits \(L_1\) and \(L_2\) with \(L_1 \neq L_2\).
    Let \(\varepsilon = \frac{\abs{L_1 - L_2}}{3} > 0\).
    By \nameref{def:convergence-of-a-sequence}, there exist \(N_1, N_2 \in \setN\) such that
    \[
      \abs{s_n - L_1} < \varepsilon \quad\text{for all } n > N_1,
    \]
    and
    \[
      \abs{s_n - L_2} < \varepsilon \quad\text{for all } n > N_2.
    \]
    Let \(N = \max(N_1, N_2)\). Then, for all \(n > N\),
    \[
      \abs{L_1 - L_2} = \abs{(L_1 - s_n) + (s_n - L_2)} \leq \abs{s_n - L_1} + \abs{s_n - L_2} < 2\varepsilon = \frac{2\abs{L_1 - L_2}}{3},
    \]
    which is a contradiction. Therefore, by contradiction, the sequence \((s_n)\) has at most one limit.
  \end{subprove}
\end{enumerate}



Let $\left(s_n\right)$ and $\left(t_n\right)$ be real-valued sequences, and $s, t \in \mathbb{R}$.
\begin{enumerate}
  \item Suppose \((s_n)\) is convergent. Prove that it is bounded. (Hint: apply the definition with \(\varepsilon = 1\).)
  \begin{define}{Bounded Sequence}[def:bounded-sequence][red]
    A sequence \((s_n)_{n\in\setN}\) in \(\setR\) is \textbf{bounded} if there exists a real number \(M>0\) such that
    \[
      |s_n| \leq M \quad\text{for all } n \in \setN.
    \]
  \end{define}
  \begin{subprove}
    Since \((s_n)\) is convergent, by \nameref{def:convergence-of-a-sequence}, there exists \(N \in \setN\) such that
    \[
      \abs{s_n - s} < 1 \quad\text{for all } n \geq N.
    \]
    Thus, for all \(n \geq N\),
    \[
      \abs{s_n} = \abs{s_n - s + s} \leq \abs{s_n - s} + \abs{s} < 1 + \abs{s}.
    \]
    Let
    \[
      M = \max\set{ \abs{s_1}, \abs{s_2}, \ldots, \abs{s_{N-1}}, 1 + \abs{s} }.
    \]
    Then, for all \(n \in \setN\), it follows that
    \[
      \abs{s_n} \leq M.
    \]
    Therefore, by \nameref{def:bounded-sequence}, the sequence \((s_n)\) is bounded.
  \end{subprove}


  \newpage
  
  \item Give an example of a nonmonotone convergent sequence other than \(((-1)^n / n)\).
  \begin{define}{Monotone Sequence}[def:monotone-sequence][red]
    A sequence \((s_n)_{n\in\setN}\) in \(\setR\) is \textbf{monotone} if it is either
    \begin{itemize}
      \item \textbf{increasing}, that is, \(s_n \leq s_{n+1}\) for all \(n \in \setN\), or
      \item \textbf{decreasing}, that is, \(s_n \geq s_{n+1}\) for all \(n \in \setN\).
    \end{itemize}
    A sequence \((s_n)_{n\in\setN}\) in \(\setR\) is \textbf{strictly monotone} if it is either
    \begin{itemize}
      \item \textbf{strictly increasing}, that is, \(s_n < s_{n+1}\) for all \(n \in \setN\), or
      \item \textbf{strictly decreasing}, that is, \(s_n > s_{n+1}\) for all \(n \in \setN\).
    \end{itemize}
  \end{define}
  \begin{subanswer}
    The sequence \(s_n = \frac{n}{n+1}\cdot (-1)^n\) is nonmonotone and convergent.
  \end{subanswer}
  
  \item Prove the Monotone Convergence Theorem.
  \begin{define}{Monotone Convergence Lemma}[lem:monotone-convergence-theorem][red]
    Every bounded monotone sequence in \(\setR\) converges.
  \end{define}
  \begin{subprove}
    Suppose \((s_n)_{n\in\setN}\) is a bounded monotone sequence in \(\setR\).
    Without loss of generality, assume \((s_n)\) is increasing and bounded above for the sake of direct proof.
    Then, by \nameref{def:completeness-axiom}, the set \(S = \set{s_n : n \in \setN}\) has a supremum \(L = \sup(S) \in \setR\).
    Let \(\varepsilon > 0\) be given.
    Since \(L - \varepsilon\) is not an upper bound for \(S\),
    there exists \(N \in \setN\) such that \(L - \varepsilon < s_N\).
    Then, for all \(n > N\), by the \nameref{def:monotone-sequence} and definition of supremum,
    \[
      L - \varepsilon < s_N \leq s_n \leq L < L + \varepsilon
    \]
    which implies
    \[
      \abs{s_n - L} < \varepsilon.
    \]
    Thus, by \nameref{def:convergence-of-a-sequence}, the sequence converges to \(L\).
  \end{subprove}
  
  \item Prove that if \((s_n)\) is an unbounded increasing sequence, 
  then \(\lim s_n = +\infty\).
  \begin{subprove}
    Let \(M > 0\) be given.
    Since \((s_n)\) is increasing, it is bounded below by \(s_1\).
    Because \((s_n)\) is unbounded, it must therefore be unbounded above.
    Hence, there must exist \(N \in \setN\) 
    such that \(s_N > M\).
    
    Since the sequence is increasing, for all \(n > N\), 
    \[
      s_n \ge s_N > M \quad\text{for all } n > N.
    \]
    Thus, by definition of divergence to infinity, \(\lim s_n = +\infty\).
  \end{subprove}
  \newpage
  \item How would you modify your previous proof to show that if \((s_n)\) is an unbounded decreasing sequence, then \(\lim s_n = -\infty\)?
  \begin{subprove}
    Let \(M > 0\) be given.
    Since \((s_n)\) is decreasing, it is bounded above by \(s_1\).
    Because \((s_n)\) is unbounded, it must therefore be unbounded below.
    Hence, there must exist \(N \in \setN\) 
    such that \(s_N < -M\).
    
    Since the sequence is decreasing, for all \(n > N\), 
    \[
      s_n \le s_N < -M \quad\text{for all } n > N.
    \]
    Thus, by definition of divergence to negative infinity, \(\lim s_n = -\infty\).
  \end{subprove}
  \item If \((s_n)\) is monotone, what can you conclude about \(\lim s_n\)? Justify your answer.
  \begin{subanswer}
    \(\lim s_n\) always exists.
    \begin{enumerate}
      \item \textbf{Case 1: \((s_n)\) is bounded.}
      By \nameref{lem:monotone-convergence-theorem}, every 
      bounded monotone sequence converges to a limit in \(\setR\).
      (Specifically, if it is increasing, it converges to 
      \(\sup\set{s_n : n \in \setN}\); if decreasing, 
      it converges to \(\inf\set{s_n : n \in \setN}\)).
      \item \textbf{Case 2: \((s_n)\) is unbounded.}
      If \((s_n)\) is increasing and unbounded, 
      then by the previous problem, \(\lim s_n = +\infty\).
      Similarly, if \((s_n)\) is decreasing and unbounded, 
      then by the previous problem, \(\lim s_n = -\infty\).
    \end{enumerate}
    Therefore, proved by cases, the sequence has a limit if \((s_n)\) is monotone.
  \end{subanswer}
  \item Let \(k \in \setR\). Prove from the definition that if the sequence \((s_n)\) 
  converges to \(s \in \setR\), then the sequence \((k s_n)\) converges to \(k s\).
  \begin{subprove}
    Suppose that \((s_n)\) converges to \(s \in \setR\).
    Let \(\varepsilon > 0\) be given.\proofstep
    \textbf{Case 1:} Suppose that \(k = 0\).
    Then, there exists \(N \in \setN\) such that
    \[
      \abs{k s_n - k s} = \abs{0 - 0} = 0 < \varepsilon \quad\text{for all } n > N.
    \]
    
    Then, there exists \(N \in \setN\) such that, by \nameref{def:convergence-of-a-sequence}, \(\lim (k s_n)= k s = 0 \).
    
    \textbf{Case 2:} Suppose that \(k \neq 0\).
    \(\frac{\varepsilon}{\abs{k}} > 0\) is also an arbitrary positive real number since \(\abs{k} > 0\).
    Then, there exists \(N \in \setN\) such that
    \[
      \abs{s_n - s} < \frac{\varepsilon}{\abs{k}} \quad\text{for all } n > N.
    \]
    Then, for all \(n > N\), it follows that
    \[
      \abs{k s_n - k s} = \abs{k} \cdot \abs{s_n - s} < \abs{k} \cdot \frac{\varepsilon}{\abs{k}} = \varepsilon.
    \]
    Thus, by definition, \(\lim (k s_n) = k s\).\proofstep
    Therefore, proved by cases, \(\lim (k s_n) = k s\).
  \end{subprove}
  \item Prove from the definition that if \((s_n)\) converges to \(s \in \setR\) and \((t_n)\) converges to \(t \in \setR\), then \((s_n + t_n)\) converges to \(s + t\).
  \begin{subprove}
    Suppose that \((s_n)\) converges to \(s \in \setR\) and \((t_n)\) converges to \(t \in \setR\).
    Let \(\varepsilon > 0\) be given. Then, 
    \(\frac{\varepsilon}{2} > 0\) is also an arbitrary positive real number.
    By the \nameref{def:convergence-of-a-sequence}, there exists \(N_1, N_2 \in \setN\) such that,
    \begin{align*}
      \abs{s_n - s} < \frac{\varepsilon}{2} \quad\text{for all } n > N_1, \\
      \abs{t_n - t} < \frac{\varepsilon}{2} \quad\text{for all } n > N_2,
    \end{align*}
    Let \(N = \max\set{N_1, N_2}\). Then, for all \(n > N\), by \nameref{def:triangle-inequality},
    \begin{align*}
      \abs{(s_n + t_n) - (s + t)} = \abs{(s_n - s) + (t_n - t)} \leq \abs{s_n - s} + \abs{t_n - t} < \frac{\varepsilon}{2} + \frac{\varepsilon}{2} = \varepsilon,
    \end{align*}
    Thus, by \nameref{def:convergence-of-a-sequence}, \(\lim (s_n + t_n) = s + t\).
  \end{subprove}
  \item Prove from the definition that if \((s_n)\) converges to \(s\) and \((t_n)\) converges to \(t\), then \((s_n t_n)\) converges to \(s t\). (Hint: use the triangle inequality, boundedness of convergent sequences, and addition of \(0 = -s_n t + s_n t\).)
  \begin{subprove}
    Suppose that \((s_n)\) converges to \(s \in \setR\) and \((t_n)\) converges to \(t \in \setR\).
    Since \((s_n)\) is convergent, by the previous result, \((s_n)\) is bounded.
    Then, by \nameref{def:bounded-sequence}, there exists \(M > 0\) such that
    \[
      \abs{s_n} \leq M \quad\text{for all } n \in \setN.
    \]
    Let \(\varepsilon > 0\) be given.
    Since \(M > 0\), it follows that \(\frac{\varepsilon}{2M} > 0\) is also an arbitrary positive real number.
    Thus, by definition, there exists \(N_1 \in \setN\) such that
    \[
      \abs{t_n - t} < \frac{\varepsilon}{2M} \quad\text{for all } n > N_1.
    \]
    Since \(\max\set{\abs{t}, 1} \geq 1 > 0\), it follows that \(\frac{\varepsilon}{2\max\set{\abs{t}, 1}} > 0\) 
    is also an arbitrary positive real number.
    Since \((s_n)\) converges to \(s\), by \nameref{def:convergence-of-a-sequence},
    there exists \(N_2 \in \setN\) such that
    \[
      \abs{s_n - s} < \frac{\varepsilon}{2\max\set{\abs{t}, 1}} \quad\text{for all } n > N_2.
    \]
    Let \(N = \max\set{N_1, N_2}\). Then, for all \(n > N\), by \nameref{def:triangle-inequality},
    \begin{align*}
      \abs{s_n t_n - s t} &= \abs{s_n t_n - s_n t + s_n t - s t} \\
      &= \abs{s_n(t_n - t) + t(s_n - s)} \\
      &\leq \abs{s_n}\abs{t_n - t} + \abs{t}\abs{s_n - s} \\
      &< M \cdot \frac{\varepsilon}{2M} + \abs{t} \cdot \frac{\varepsilon}{2\max\set{\abs{t}, 1}} \\
      &= \frac{\varepsilon}{2} + \frac{\abs{t}}{\max\set{\abs{t}, 1}} \cdot \frac{\varepsilon}{2} \\
      &\leq \frac{\varepsilon}{2} + \frac{\varepsilon}{2} = \varepsilon,
    \end{align*}
    Thus, by \nameref{def:convergence-of-a-sequence}, \(\lim (s_n t_n) = s t\).
  \end{subprove}











  \newpage
  \item Prove the following Cauchy Criterion (forward direction).
  \begin{define}{Cauchy Convergence Criterion}[thm:cauchy-convergence-criterion][red]
    A sequence \((s_n)_{n\in\setN}\) in \(\setR\) converges if and only if it is a Cauchy sequence.
  \end{define}
  \begin{subprove}
    (\(\Rightarrow\)) Suppose that \((s_n)\) converges to a limit \(s \in \setR\).
    Let \(\varepsilon > 0\) be given. Then, \(\dfrac{\varepsilon}{2} > 0\) is also an arbitrary positive real number.
    By \nameref{def:convergence-of-a-sequence}, there exists \(N \in \setN\) such that
    \[
      \abs{s_k - s} < \dfrac{\varepsilon}{2} \quad\text{for all } k > N.
    \]
    Let \(n, m > N\) be arbitrary integers. Then, by the \nameref{def:triangle-inequality},
    \[
      \abs{s_n - s_m} \leq \abs{s_n - s} + \abs{s_m - s} < \dfrac{\varepsilon}{2} + \dfrac{\varepsilon}{2} = \varepsilon.
    \]
    Therefore, by \nameref{def:cauchy-sequence}, \((s_n)\) is a Cauchy sequence.

    % (\(\Leftarrow\)) Suppose that \((s_n)\) is a Cauchy sequence.
    % Let \(\varepsilon > 0\) be given. Then \(\dfrac{\varepsilon}{2} > 0\).
    % By \nameref{def:cauchy-sequence}, there exists \(N \in \setN\) such that
    % \[
    %   \abs{s_n - s_m} < \dfrac{\varepsilon}{2} \quad\text{for all } n, m > N.
    % \]
    % Let \(s = \lim s_n\). Then, by \nameref{def:convergence-of-a-sequence}, there exists \(N \in \setN\) such that
    % \[
    %   \abs{s_n - s} < \dfrac{\varepsilon}{2} \quad\text{for all } n > N.
    % \]
    % Then, for all \(n, m > N\), by the \nameref{def:triangle-inequality},
    % \[
    %   \abs{s_n - s_m} \leq \abs{s_n - s} + \abs{s_m - s} < \dfrac{\varepsilon}{2} + \dfrac{\varepsilon}{2} = \varepsilon.
    % \]
    % Therefore, by \nameref{def:convergence-of-a-sequence}, \((s_n)\) converges to \(s\).
  \end{subprove}

  % \item Prove the following Cauchy Criterion (backward direction).
  % \begin{define}{Bolzano-Weierstrass Theorem}[thm:bolzano-weierstrass][red]
  %   Every bounded sequence of real numbers has a convergent subsequence.
  % \end{define}
  % \begin{subprove}
  %   Suppose that \((s_n)\) is a Cauchy sequence.
  %   First, as proven previously, every Cauchy sequence is bounded.
  %   Since \((s_n)\) is bounded, by the \nameref{thm:bolzano-weierstrass}, there exists a subsequence \((s_{n_k})\) that converges to some limit \(s \in \setR\).
  
  %   We claim that the original sequence \((s_n)\) converges to \(s\).
  %   Let \(\varepsilon > 0\) be given. Then \(\frac{\varepsilon}{2} > 0\).
    
  %   Since \((s_n)\) is Cauchy, by \nameref{def:cauchy-sequence}, there exists \(N_1 \in \setN\) such that
  %   \[
  %     \abs{s_n - s_m} < \frac{\varepsilon}{2} \quad\text{for all } n, m > N_1.
  %   \]
  %   Since the subsequence \((s_{n_k})\) converges to \(s\), there exists \(K \in \setN\) such that
  %   \[
  %     \abs{s_{n_k} - s} < \frac{\varepsilon}{2} \quad\text{for all } k > K.
  %   \]
  %   Let \(N = \max\set{N_1, K}\). Let \(n > N\).
  %   We choose a specific index \(n_k\) from the subsequence such that \(k > N\). 
  %   Note that \(n_k \ge k > N \ge N_1\).
    
  %   Then, by the \nameref{def:triangle-inequality}:
  %   \begin{align*}
  %     \abs{s_n - s} 
  %     &= \abs{(s_n - s_{n_k}) + (s_{n_k} - s)} \\
  %     &\leq \abs{s_n - s_{n_k}} + \abs{s_{n_k} - s}.
  %   \end{align*}
  %   Since \(n > N_1\) and \(n_k > N_1\), the first term \(\abs{s_n - s_{n_k}} < \frac{\varepsilon}{2}\) by the Cauchy property.
  %   Since \(k > K\), the second term \(\abs{s_{n_k} - s} < \frac{\varepsilon}{2}\) by the subsequence convergence.
    
  %   Therefore,
  %   \[
  %     \abs{s_n - s} < \frac{\varepsilon}{2} + \frac{\varepsilon}{2} = \varepsilon.
  %   \]
  %   Thus, by \nameref{def:convergence-of-a-sequence}, \((s_n)\) converges to \(s\).
  % \end{subprove}


  
  % \item Prove the Bolzano-Weierstrass Theorem: Every bounded sequence of real numbers has a convergent subsequence.
  % \begin{define}{Nested Interval Property}[thm:nested-interval-property][red]
  %   Let \(I_1 \supseteq I_2 \supseteq I_3 \supseteq \dots\) be a sequence of nested closed bounded intervals 
  %   \(I_k = [a_k, b_k]\). If the lengths of the intervals converge to 0 (i.e., \(\lim (b_k - a_k) = 0\)), 
  %   then their intersection contains exactly one point: \(\bigcap_{k=1}^{\infty} I_k = \set{x}\).
  % \end{define}
  % \begin{subprove}
  %   Let \((s_n)\) be a bounded sequence. 
  %   Since it is bounded, there exist \(a, b \in \setR\) such that \(s_n \in [a, b]\) for all \(n \in \setN\).
  %   Let \(I_1 = [a, b]\).
    
  %   We construct a sequence of nested intervals recursively:
  %   \begin{itemize}
  %     \item Divide \(I_1\) into two equal sub-intervals: \([a, \frac{a+b}{2}]\) and \([\frac{a+b}{2}, b]\).
  %     \item Since \(I_1\) contains the infinite set of sequence terms \(\set{s_n : n \in \setN}\), at least one of these two sub-intervals must contain infinitely many terms of the sequence. (If both had finitely many, their union would be finite).
  %     \item Let \(I_2\) be the half that contains infinitely many terms.
  %     \item Repeat this process. Given \(I_k\), divide it in half. Choose \(I_{k+1}\) to be the half containing infinitely many terms of the sequence.
  %   \end{itemize}
    
  %   This generates a sequence of nested intervals \(I_1 \supseteq I_2 \supseteq \dots\) with lengths:
  %   \[
  %     \text{length}(I_k) = \frac{b-a}{2^{k-1}}.
  %   \]
  %   Since \(\lim_{k\to\infty} \frac{b-a}{2^{k-1}} = 0\), by the \nameref{thm:nested-interval-property}, 
  %   there exists a unique point \(s \in \bigcap_{k=1}^{\infty} I_k\).
  
  %   Now, we construct the convergent subsequence \((s_{n_k})\):
  %   \begin{itemize}
  %     \item Choose \(n_1\) such that \(s_{n_1} \in I_1\).
  %     \item Since \(I_2\) contains infinitely many terms, choose \(n_2 > n_1\) such that \(s_{n_2} \in I_2\).
  %     \item generally, since \(I_k\) contains infinitely many terms, we can always choose \(n_k > n_{k-1}\) such that \(s_{n_k} \in I_k\).
  %   \end{itemize}
    
  %   For any \(k\), both \(s_{n_k}\) and \(s\) are in the interval \(I_k\). Therefore, the distance between them is bounded by the length of the interval:
  %   \[
  %     \abs{s_{n_k} - s} \leq \text{length}(I_k) = \frac{b-a}{2^{k-1}}.
  %   \]
  %   As \(k \to \infty\), \(\frac{b-a}{2^{k-1}} \to 0\). 
  %   Thus, by the Squeeze Theorem (or definition of convergence), \(\lim_{k\to\infty} s_{n_k} = s\).
  %   Therefore, \((s_n)\) has a convergent subsequence.
  % \end{subprove}

  \item Prove that the sequence \(s_n=(-1)^{n+1}\dfrac{n}{n+1}\) is not Cauchy. 
  Conclude it is divergent.
  \begin{define}{Cauchy Sequence}[def:cauchy-sequence][red]
    A sequence \((s_n)_{n\in\setN}\) in \(\setR\) is a \textbf{Cauchy sequence} if,
    for every \(\varepsilon>0\), there exists \(N\in\setN\) such that
    \[
      \abs{s_n - s_m} < \varepsilon \quad\text{for all } n, m > N.
    \]
  \end{define}  
  \begin{subprove}
    Suppose that the sequence \(s_n=(-1)^{n+1}\dfrac{n}{n+1}\) is Cauchy 
    for the sake of contradiction.
    Then, by \nameref{def:cauchy-sequence}, for all \(\varepsilon > 0\), there exists \(N \in \setN\) such that
    \[
      \abs{s_n - s_m} < \varepsilon \quad\text{for all } n, m > N.
    \]
    Let \(\varepsilon = 0.1\) and \(n=m+1\). Then,
    \[
      0.5 =\abs{2 - \frac{1}{1} - \frac{1}{2}} \le \abs{2 - \frac{1}{m+1} - \frac{1}{m+2}} = \abs{\dfrac{m+1}{m+2} + \dfrac{m}{m+1}} = \abs{s_n - s_m} < 0.1.
    \]
    which is a contradiction.
    Thus, by contradiction, the sequence \(s_n=(-1)^{n+1}\dfrac{n}{n+1}\) is not Cauchy.
    Therefore, by the contrapositive of \nameref{thm:cauchy-convergence-criterion}, the sequence is divergent.
  \end{subprove}


  \newpage
  \item Prove that every real-valued Cauchy sequence is bounded.
  \begin{subprove}
    Suppose that \((s_n)_{n\in\setN}\) is a Cauchy sequence in \(\setR\) for the sake of direct proof.
    By \nameref{def:cauchy-sequence}, there exists \(N \in \setN\) such that
    \[
      \abs{s_n - s_m} < 1 \quad\text{for all } n, m > N.
    \]
    Let \(m = N+1\). Then, for all \(n > N\),
    \[
      \abs{s_n - s_{N+1}} < 1.
    \]
    Thus, for all \(n > N\), by \nameref{def:triangle-inequality},
    \[
      \abs{s_n} = \abs{s_n - s_{N+1} + s_{N+1}} \leq \abs{s_n - s_{N+1}} + \abs{s_{N+1}} < 1 + \abs{s_{N+1}}.
    \]
    Let \(M = \max\set{ \abs{s_1}, \abs{s_2}, \ldots, \abs{s_N}, 1 + \abs{s_{N+1}} }\). Then, for all \(n \in \setN\), it follows that
    \[
      \abs{s_n} \leq M.
    \]
    Therefore, by \nameref{def:bounded-sequence}, the sequence \((s_n)\) is bounded.
  \end{subprove}

  \item There is a theorem: every real-valued Cauchy sequence converges to a real number. 
  Is it true that every rational-valued Cauchy sequence converges to a rational number? Justify your answer.
  \begin{subanswer}
    No, it is not true that every rational-valued Cauchy sequence converges to a rational number.
    Consider the sequence \((s_n)_{n\in\setN}\) defined by
    \[
      s_n = \frac{\floor{\sqrt{2}n}}{n}.
    \]
    Then, for all \(n\in \setN\), by the property of the floor function,
    \[
      \dfrac{\sqrt{2}n - 1}{n} \le s_n \le \dfrac{\sqrt{2}n}{n} = \sqrt{2}.
    \]
    Thus, by the \nameref{thm:squeeze-theorem}, the sequence \((s_n)\) converges to \(\sqrt{2}\).
    Since \(\sqrt{2}\) is irrational, the sequence \((s_n)\) does not converge to a rational number.
    Therefore, by counterexample, not every rational-valued Cauchy sequence converges to a rational number.
  \end{subanswer}



  \newpage
  \item Suppose \((s_n)\) converges.
    \begin{enumerate}
      \item Prove from the definition that if \(s_n\ge a\) 
      for all but finitely many \(n\), then \(\lim s_n\ge a\).
      \begin{subprove}
        Let \(s = \lim s_n\). Suppose for the sake of contradiction that \(s < a\).
        Let \(\varepsilon = a - s\). Since \(s < a\), it follows that \(\varepsilon > 0\).
        By \nameref{def:convergence-of-a-sequence}, there exists \(N_1 \in \setN\) such that
        \[
          \abs{s_n - s} < \varepsilon \quad\text{for all } n > N_1.
        \]
        In particular, this implies \(s_n - s < \varepsilon\), so
        \[
          s_n < s + \varepsilon = s + (a - s) = a.
        \]
        Thus, \(s_n < a\) for all \(n > N_1\).
        However, by the assumption,
        there exists \(N_2 \in \setN\) such that \(s_n \ge a\) for all \(n > N_2\).
        
        Let \(n > \max\set{N_1, N_2}\). Then, \(a \le s_n < a\)
        which is a contradiction. Therefore, \(\lim s_n \ge a\).
      \end{subprove}
      \item If all but finitely many \(s_n\) belong to \([a,b]\), 
      what can you conclude about \(\lim s_n\)? Justify your answer.\begin{subanswer}
        Conclusion: \(\lim s_n \in [a, b]\).
    
        Justification:
        Let \(s = \lim s_n\).
        The condition \(s_n \in [a, b]\) for all but finitely many \(n\) is equivalent to saying:
        \[
          s_n \ge a \quad\text{and}\quad s_n \le b \quad\text{for all but finitely many } n.
        \]
        Using the result proven in part (a):
        \begin{itemize}
          \item Since \(s_n \ge a\) eventually, it follows that \(s \ge a\).
          \item Similarly, since \(s_n \le b\) implies \(-s_n \ge -b\), applying the same logic shows \(-s \ge -b\), or \(s \le b\).
        \end{itemize}
        Combining these inequalities, we have \(a \le s \le b\). 
        Thus, the limit lies in the closed interval \([a, b]\).
      \end{subanswer}
    \end{enumerate}
  \item Suppose there exists \(N_0\in\mathbb{N}\) such that \(s_n\le t_n\) for all \(n>N_0\).
    \begin{enumerate}
      \item Prove that if \(\lim s_n=+\infty\), then \(\lim t_n=+\infty\).
      \begin{subprove}
        Let \(M > 0\) be given. Then, by the \nameref{def:divergence-to-infinity}, 
        there exists \(N_1 \in \setN\) such that
        \[
          s_n > M \quad\text{for all } n > N_1.
        \]
        Let \(N = \max\set{N_0, N_1}\). Then, for all \(n > N\), by the assumption,
        \[
          t_n \ge s_n > M.
        \]
        Therefore, by definition, \(\lim t_n = +\infty\).
      \end{subprove}
      \newpage
      \item Prove that if \(\lim t_n=-\infty\), then \(\lim s_n=-\infty\).
      \begin{subprove}
        Let \(M < 0\) be given. Then, by the \nameref{def:divergence-to-infinity}, 
        there exists \(N_1 \in \setN\) such that
        \[
          t_n < M \quad\text{for all } n > N_1.
        \]
        Let \(N = \max\set{N_0, N_1}\). Then, for all \(n > N\), by the assumption,
        \[
          s_n \le t_n < M.
        \]
        Therefore, by definition, \(\lim s_n = -\infty\).
      \end{subprove}
      \item Prove that if \(\lim s_n\) and \(\lim t_n\) exist (finite), then \(\lim s_n\le \lim t_n\).
      \begin{subprove}
        Suppose that \(\lim s_n = s\) and \(\lim t_n = t\) and \(s - t > 0\) for the sake of contradiction.
        Let \(\varepsilon > 0\) be given. Then, \(\dfrac{\varepsilon}{2} > 0\) is also an arbitrary positive real number.
        By \nameref{def:convergence-of-a-sequence}, there exists \(N_1 \in \setN\) such that
        \[
          \abs{s_n - s} < \dfrac{\varepsilon}{2} \quad\text{for all } n > N_1
        \]
        and there exists \(N_2 \in \setN\) such that
        \[
          \abs{t_n - t} < \dfrac{\varepsilon}{2} \quad\text{for all } n > N_2.
        \]
        Let \(N = \max\set{N_1, N_2}\). Then, for all \(n > N\), by the assumption,
        \[
          s - \dfrac{\varepsilon}{2} < s_n \le t_n < t + \dfrac{\varepsilon}{2}.
        \]
        That is, with \(\varepsilon = s - t\), we have
        \[
          s - t < s - t.
        \]
        This is a contradiction. Therefore, by contradiction, \(\lim s_n\le \lim t_n\).
      \end{subprove}
    \end{enumerate}
\end{enumerate}



























%%%%%%%%%%%%%%%%%%%%%%%%%%%%%%%%%%%%%%%%%%%%%%%%%%%%%%%%%%%%%%%
% Analysis 2 - HW
%%%%%%%%%%%%%%%%%%%%%%%%%%%%%%%%%%%%%%%%%%%%%%%%%%%%%%%%%%%%%%%
\newpage\handout[true]
{\textsc{Math 2106-B: Foundations of Mathematical Proof}}{Fall 2025}
{Homework 9: Sequences and Limits}
{\textit{Prof.} \href{https://ceheitsch.github.io/webpage/}{Dr. Christine Heitsch}}{Due Date: December 1, 2025}
{Math 2106-B: HW9}
{\large
  \noindent
  \textbf{Name:} \HandoutAuthor \smallskip \\

  \noindent
  \textbf{GTID:} 903743506 \smallskip \\

  \noindent
  %%% List anyone with whom you discussed the problems
  \textbf{Collaborators:} None. This material reflects independent preparation. \smallskip \\

  \noindent
  %%% List all resources used OTHER than the textbook, lecture notes, 
  %%% webwork problems, ICA questions, and previous homework assignments.
  \textbf{Outside resources:} This work was informed by MATH 2106 course materials 
  (ICAs: in-class activities and WebWork contents) from \HandoutInstitution\ (the course textbook is 
  Chartrand \parencite{chartrand_mathematical_2018}).   
  Outside references specifically include
  textbooks written by 
  Grimaldi \parencite{grimaldi_discrete_2004} (used in MATH 3012) and 
  Rosen \parencite{rosen_discrete_2019} (used in CS 2050). \smallskip \\

  \noindent
  \textbf{Notice}: The answers contain cross-references throughout this document as clickable hyperlinks in most PDF viewers. However, the Gradescope PDF view might not support clickable hyperlinks.
}
\printbibliography[
  heading=bibintoc,   
  title={References}     % change “References” text if needed
]
% `heading=bibnumbered` if you prefer it numbered chapter (in \tableofcontents)
% `heading=bibintoc`    if you prefer it unnumbered but still in the ToC
\newpage
%%%%%%%%%%%%%%%%%%%%%%%%%%%%%%%%%%%%%%%%%%%%%%%%%%%%%%%%%%%%%%%%%%%%%%%
% Problems Section
%%%%%%%%%%%%%%%%%%%%%%%%%%%%%%%%%%%%%%%%%%%%%%%%%%%%%%%%%%%%%%%%%%%%%%%
Let \((s_n)\) and \((t_n)\) be real-valued sequences.
\begin{enumerate}
  \item 
  \begin{enumerate}
    \item Prove that \(\lim s_n = 0\) if and only if \(\lim |s_n| = 0\).
    \begin{subprove}
      (\(\Rightarrow\)) Suppose that \(\lim s_n = 0\) for the sake of direct proof.
      Let \(\varepsilon > 0\) be given. Then, by \nameref{def:convergence-of-a-sequence},
      there exists \(N \in \setN\) such that 
      \[
        \abs{s_n - 0} = \abs{s_n} < \varepsilon \quad\text{for all } n > N.
      \]
      Since \(\abs{\abs{s_n} - 0} = \abs{\abs{s_n}} = \abs{s_n} = \abs{s_n - 0}\), it follows that
      \[
        \abs{\abs{s_n} - 0} < \varepsilon \quad\text{for all } n > N.
      \]
      Therefore, by definition, \(\lim |s_n| = 0\).

      (\(\Leftarrow\)) Suppose that \(\lim |s_n| = 0\) for the sake of direct proof.
      Let \(\varepsilon > 0\) be given. Then, by definition,
      there exists \(N \in \setN\) such that 
      \[
        \abs{|s_n| - 0} = |s_n| < \varepsilon \quad\text{for all } n > N.
      \]
      Since \(\abs{s_n - 0} = \abs{s_n} = \abs{\abs{s_n}} = \abs{\abs{s_n}-0}\), it follows that
      \[
        \abs{s_n - 0} < \varepsilon \quad\text{for all } n > N.
      \]
      Therefore, by definition, \(\lim s_n = 0\).
    \end{subprove}
    \item Give an example where \(\lim |s_n|\) exists but \(\lim s_n\) does not. 
    Justify your answer.
    \begin{subanswer}
      Let \(s_n = (-1)^n\) for all \(n \in \setN\).
      Since \(\abs{s_n} = \abs{(-1)^n} = 1\) for all \(n \in \setN\), the sequence \((\abs{s_n})\) is constant.
      Let \(\varepsilon > 0\) be given. Then, 
      for any \(n, N \in \setN\) where \(n > N\), 
      \[
        \abs{\abs{s_n} - 1} = \abs{1 - 1} = 0 < \varepsilon \quad\text{for all } n > N.
      \]
      Therefore, by definition, \(\lim \abs{s_n} = 1\).
      
      Next, suppose for the sake of contradiction that \(\lim s_n\) exists.
      Then, by definition, there exists \(s \in \setR\). 
      Let \(\varepsilon = 1\). There exists \(N \in \setN\) such that 
      \[
        \abs{s_n - s} < 1 \quad\text{for all } n > N.
      \]
      Now, choose \(k \in \setN\) such that \(2k > N\).
      Since \(s_n\) alternates, \(s_{2k} = 1\) and \(s_{2k+1} = -1\).
      Since \(2k > N\) and \(2k+1 > N\), it follows that
      \[
        2 = \abs{1 - (-1)} = \abs{s_{2k} - s_{2k+1}} \leq \abs{s_{2k} - s} + \abs{s - s_{2k+1}} < 1 + 1 = 2.
      \]
      This implies \(2 < 2\), which is a contradiction. 
      Therefore, \(\lim s_n\) does not exist.
    \end{subanswer}
  \end{enumerate}



  \newpage

  \item Suppose \(\lim t_n = 0\).
  \begin{enumerate}
    \item Prove that if \((s_n)\) is bounded, then \(\lim (s_n t_n) = 0\).
    \begin{subprove}
      Suppose that \((s_n)\) is bounded and \(\lim t_n = 0\) for the sake of direct proof.
      Since \((s_n)\) is bounded, by \nameref{def:bounded-sequence}, there exists \(M > 0\) such that
      \[
        \abs{s_n} \leq M \quad\text{for all } n \in \setN.
      \]
      Let \(\varepsilon > 0\) be given. Since \(M > 0\), it follows that \(\varepsilon / M > 0\) 
      is also an arbitrary positive real number.
      Then, by definition, since \(\lim t_n = 0\),
      there exists \(N \in \setN\) such that 
      \[
        \abs{t_n - 0} = \abs{t_n} < \frac{\varepsilon}{M} \quad\text{for all } n > N.
      \]
      Then, for all \(n > N\), it follows that
      \[
        \abs{s_n t_n - 0} = \abs{s_n t_n} = \abs{s_n} \cdot \abs{t_n} \leq M \cdot \abs{t_n} < M \cdot \frac{\varepsilon}{M} = \varepsilon.
      \]
      Therefore, by definition, \(\lim (s_n t_n) = 0\).
    \end{subprove}
    \item Give an example where \(\lim (s_n t_n)\) is not zero. Justify your answer.
    \begin{subanswer}
      Consider the sequences defined by \(s_n = n\) and \(t_n = 1/n\) for all \(n \in \setN\)
      since it is trivial that \(\lim t_n = \lim 1/n = 0\). Then, 
      \[
        s_n t_n = n \cdot \frac{1}{n} = 1 \quad\text{for all } n \in \setN.
      \]
      The sequence \((s_n t_n)\) is the constant sequence \((1, 1, 1, \dots)\).
      Clearly,
      \[
        \lim (s_n t_n) = 1 \neq 0.
      \]
      This is possible because the sequence \((s_n)\) is not bounded.
    \end{subanswer}
  \end{enumerate}


  \newpage

  \item Prove that if \((s_n)\) is a sequence of positive real numbers, 
  then \(\lim s_n = +\infty\) if and only if \(\lim (1 / s_n) = 0\).
  \begin{subprove}
    (\(\Rightarrow\)) Suppose that \(\lim s_n = +\infty\) for the sake of direct proof.
    Let \(\varepsilon > 0\) be given. Since \(\varepsilon > 0\), it follows that \(\frac{1}{\varepsilon} > 0\) 
    is also an arbitrary positive real number.
    Then, by \nameref{def:divergence-to-infinity},
    there exists \(N \in \setN\) such that
    \[
      s_n > \frac{1}{\varepsilon} \quad\text{for all } n > N.
    \]
    Since \(s_n > 0\) for all \(n \in \setN\), it follows that
    \[
      \abs{\frac{1}{s_n} - 0} = \frac{1}{s_n} < \varepsilon \quad\text{for all } n > N.
    \]
    Therefore, by definition, \(\lim (1 / s_n) = 0\).

    (\(\Leftarrow\)) Suppose that \(\lim (1 / s_n) = 0\) for the sake of direct proof.
    Let \(M > 0\) be given. Since \(M > 0\), it follows that \(\frac{1}{M} > 0\) 
    is also an arbitrary positive real number.
    Then, by definition,
    there exists \(N \in \setN\) such that
    \[
      \abs{\frac{1}{s_n} - 0} = \frac{1}{s_n} < \frac{1}{M} \quad\text{for all } n > N.
    \]
    Since \(s_n > 0\) for all \(n \in \setN\), it follows that
    \[
      s_n > M \quad\text{for all } n > N.
    \]
    Therefore, by definition, \(\lim s_n = +\infty\).
  \end{subprove}



  \item Let \(S\) be a bounded nonempty subset of \(\setR\) and 
  suppose \(\sup S \notin S\). Prove that there is an increasing 
  sequence \((s_n)\) of points in \(S\) such that \(\lim s_n = \sup S\).
  \begin{subprove}
    Let \(L = \sup S\). Since \(L \notin S\) and \(L\) is the least upper bound of \(S\),
    it follows that for any \(\varepsilon > 0\), there exists \(s \in S\) such that \(L - \varepsilon < s < L\).
    Otherwise, if all \(s \in S\) satisfy \(s \leq L - \varepsilon\), then \(L - \varepsilon\) would be
    an upper bound of \(S\) smaller than \(L\), contradicting that \(L\) is the least upper bound.

    A sequence \((s_n)\) in \(S\) is recursively defined as follows.
    For \(n = 1\), let \(s_1 \in S\) such that \(L - 1 < s_1 < L\).
    For \(n \geq 2\), let \(s_n \in S\) such that
    \[
      \max\set{s_{n-1}, L - \frac{1}{n}} < s_n < L.
    \]
    Such a choice is possible because \(s_{n-1} < L\) implies that \(\max\{s_{n-1}, L - 1/n\} < L\).
    Since \(L\) is the least upper bound, there exists \(s \in S\) strictly between this maximum and \(L\).
    Also, by the recursive definition, \((s_n)\) is strictly increasing.

    Let \(\varepsilon > 0\) be given and \(N \in \setN\) be \(\ceil{1/\varepsilon}\) where \(N \ge 1/\varepsilon\).
    Then, for all \(n > N\), it follows that
      \[
        s_n > L - \frac{1}{n} > L - \frac{1}{N} \ge L - \varepsilon.
      \]
      Since \(s_n < L\) for all \(n \in \setN\) by definition, it follows that
      \[
        \abs{s_n - L} = L - s_n < \varepsilon \quad\text{for all } n > N.
      \]
    Thus, by definition, \(\lim s_n = L = \sup S\). Therefore, there is an increasing 
    sequence \((s_n)\) of points in \(S\) such that \(\lim s_n = \sup S\).
  \end{subprove}



  \newpage
  \item 
  \begin{enumerate}
    \item Prove the following Squeeze Theorem.
    \begin{define}{Squeeze Theorem}[thm:squeeze-theorem][red]
      Let \((a_n)_{n\in\setN}\), \((s_n)_{n\in\setN}\), and \((b_n)_{n\in\setN}\) be sequences in \(\setR\).
      If \(a_n \leq s_n \leq b_n\) for all \(n \in \setN\) and \(\lim a_n = \lim b_n = s \in \setR\),
      then \(\lim s_n = s\).
    \end{define}
    \textit{Hint: show that \(s - \varepsilon < s_n < s + \varepsilon\) for \(n\) sufficiently large.}
    \begin{subprove}
      Suppose that \(a_n \leq s_n \leq b_n\) for all \(n \in \setN\) and \(\lim a_n = \lim b_n = s\) for the sake of direct proof.
      Let \(\varepsilon > 0\) be given. Then, by definition, since \(\lim a_n = s\),
      there exists \(N_1 \in \setN\) such that
      \[
        \abs{a_n - s} < \varepsilon \quad\text{for all } n > N_1.
      \]
      That is,
      \[
        s - \varepsilon < a_n < s + \varepsilon \quad\text{for all } n > N_1.
      \]
      Similarly, since \(\lim b_n = s\), by definition,
      there exists \(N_2 \in \setN\) such that
      \[
        \abs{b_n - s} < \varepsilon \quad\text{for all } n > N_2.
      \]
      That is,
      \[
        s - \varepsilon < b_n < s + \varepsilon \quad\text{for all } n > N_2.
      \]
      Let \(N = \max\{N_1, N_2\}\). Then, for all \(n > N\), it follows that
      \[
        s - \varepsilon < a_n \leq s_n \leq b_n < s + \varepsilon.
      \]
      Thus,
      \[
        s - \varepsilon < s_n < s + \varepsilon \quad\text{for all } n > N,
      \]
      which implies that
      \[
        \abs{s_n - s} < \varepsilon \quad\text{for all } n > N.
      \]
      Therefore, by \nameref{def:convergence-of-a-sequence}, \(\lim s_n = s\).
      Hence, \nameref{thm:squeeze-theorem} is proved.
    \end{subprove}
    \item Suppose that \(\abs{s_n} \leq t_n\) for all \(n\) and \(\lim t_n = 0\). 
    Prove that \(\lim s_n = 0\).
    \begin{subprove}
      Suppose that \(\abs{s_n} \leq t_n\) for all \(n \in \setN\) and \(\lim t_n = 0\) for the sake of direct proof.
      Since \(0 \leq \abs{s_n} \leq t_n\), the terms \(t_n\) are non-negative for all \(n\).
      Let \(\varepsilon > 0\) be given. Then, by definition, since \(\lim t_n = 0\),
      there exists \(N \in \setN\) such that
      \[
        \abs{t_n - 0} = \abs{t_n} = t_n < \varepsilon \quad\text{for all } n > N.
      \]
      Since \(\abs{s_n} \leq t_n\) for all \(n \in \setN\), it follows that
      \[
        \abs{s_n - 0} = \abs{s_n} \leq t_n < \varepsilon \quad\text{for all } n > N.
      \]
      Therefore, by definition, \(\lim s_n = 0\).
    \end{subprove}
  \end{enumerate}



  \newpage
  \item \label{prob:limit-inequality} Let \((s_n)\) be a convergent sequence and suppose that 
  \(\lim s_n > a\). Prove that there exists a number \(N\) such that 
  \(n > N\) implies \(s_n > a\).
  \begin{subprove}
    Suppose that \((s_n)\) is a convergent sequence and \(\lim s_n = L\) where \(L > a\)
    for the sake of direct proof.
    Let \(\varepsilon = L - a\). Since \(L > a\), it follows that \(\varepsilon > 0\).
    Then, by definition, since \(\lim s_n = L\),
    there exists \(N \in \setN\) such that
    \[
      \abs{s_n - L} < \varepsilon \quad\text{for all } n > N.
    \]
    This implies that
    \[
      L - \varepsilon < s_n < L + \varepsilon \quad\text{for all } n > N.
    \]
    Substituting \(\varepsilon = L - a\) into the lower bound yields
    \[
      s_n > L - (L - a) = a \quad\text{for all } n > N.
    \]
    Therefore, there exists a number \(N\) such that \(n > N\) implies \(s_n > a\).
  \end{subprove}

  \item Let \((s_n)\) and \((t_n)\) be sequences such that 
  \(\lim s_n = +\infty\) and \(\lim t_n > 0\) (note that \((t_n)\) 
  could diverge to \(+\infty\) also). Prove from the definition 
  that \(\lim s_n t_n = +\infty\). \\
  \textit{Hint: do you know or can you prove that there 
  exists a number \(m\) such that \(0 < m < \lim t_n\)?}\begin{subprove}
    Suppose that \(\lim s_n = +\infty\) and \(\lim t_n > 0\). 
    
    First, it is established that there exists a constant \(m > 0\) and an integer \(N_1\) such that \(t_n > m\) for all \(n > N_1\). 
    The limit of \(t_n\) will be considered in two mutually exclusive cases as follows.
    
    \begin{enumerate}
      \item Suppose \(\lim t_n = L\) where \(L \in \setR\) and \(L > 0\).
      Let \(m = L/2\). Since \(L > 0\), it follows that \(m > 0\).
      Since \(\lim t_n = L > m\), by question (\ref{prob:limit-inequality}) (with \(a=m\)), there exists \(N_1 \in \setN\) such that 
      \[
        t_n > m \quad\text{for all } n > N_1.
      \]

      \item Suppose \(\lim t_n = +\infty\).
      By \nameref{def:divergence-to-infinity}, there exists \(N_1 \in \setN\) such that 
      \[
        t_n > 1 \quad\text{for all } n > N_1.
      \]
    \end{enumerate}

    Thus, a constant \(m > 0\) and an index \(N_1\) have been found such that \(t_n > m\) for all \(n > N_1\).
    Furthermore, by the assumption, \(\lim s_n = +\infty\).
    Let \(K > 0\) be an arbitrary positive real number.
    Since \(m > 0\), it follows that \(\frac{K}{m} > 0\) is also an arbitrary positive real number.
    Then, by definition, there exists \(N_2 \in \setN\) such that
    \[
      s_n > \frac{K}{m} > 0 \quad\text{for all } n > N_2.
    \]
    Let \(N = \max\{N_1, N_2\}\). Then, for all \(n > N\), it follows that
    \[
      s_n t_n > \frac{K}{m} \cdot m = K > 0.
    \]
    Since \(K\) was arbitrary, by definition, \(\lim s_n t_n = +\infty\).
  \end{subprove}








  \newpage

  \item \textbf{Complete but do not hand in}. 
  Let \(k \in \setR\). Prove from the definition that if the 
  sequence \((s_n)\) converges to \(s \in \setR\), then the 
  sequence \((k s_n)\) converges to \(k s\).
  \begin{subprove}
    Suppose that \((s_n)\) converges to \(s \in \setR\) for the sake of direct proof.
    Two cases are considered.
    
    \textbf{Case 1:} Suppose that \(k = 0\).
    Then, \(k s_n = 0\) for all \(n \in \setN\), so the sequence \((k s_n)\) is constant.
    Let \(\varepsilon > 0\) be given. Then, for any \(N \in \setN\), it follows that
    \[
      \abs{k s_n - k s} = \abs{0 - 0} = 0 < \varepsilon \quad\text{for all } n \geq N.
    \]
    Therefore, by definition, \(\lim (k s_n) = 0 = k s\).
    
    \textbf{Case 2:} Suppose that \(k \neq 0\).
    Let \(\varepsilon > 0\) be given. Since \(\abs{k} > 0\), it follows that \(\tfrac{\varepsilon}{\abs{k}} > 0\) 
    is also an arbitrary positive real number.
    Then, by definition, since \(\lim s_n = s\),
    there exists \(N \in \setN\) such that
    \[
      \abs{s_n - s} < \frac{\varepsilon}{\abs{k}} \quad\text{for all } n \geq N.
    \]
    Then, for all \(n \geq N\), it follows that
    \[
      \abs{k s_n - k s} = \abs{k} \cdot \abs{s_n - s} < \abs{k} \cdot \frac{\varepsilon}{\abs{k}} = \varepsilon.
    \]
    Therefore, by definition, \(\lim (k s_n) = k s\).
    
    In both cases, it is shown that \(\lim (k s_n) = k s\).
  \end{subprove}
  \item \textbf{Complete but do not hand in}. Prove from the 
  definition that if \((s_n)\) converges to \(s \in \setR\) 
  and \((t_n)\) converges to \(t \in \setR\), then 
  \((s_n + t_n)\) converges to \(s + t\).
  \begin{subprove}
    Suppose that \((s_n)\) converges to \(s \in \setR\) and \((t_n)\) converges to \(t \in \setR\) 
    for the sake of direct proof.
    Let \(\varepsilon > 0\) be given. Since \(\varepsilon > 0\), it follows that \(\varepsilon / 2 > 0\) 
    is also an arbitrary positive real number.
    Then, by definition, since \(\lim s_n = s\),
    there exists \(N_1 \in \setN\) such that
    \[
      \abs{s_n - s} < \frac{\varepsilon}{2} \quad\text{for all } n \geq N_1.
    \]
    Similarly, since \(\lim t_n = t\), by definition,
    there exists \(N_2 \in \setN\) such that
    \[
      \abs{t_n - t} < \frac{\varepsilon}{2} \quad\text{for all } n \geq N_2.
    \]
    Let \(N = \max\{N_1, N_2\}\). Then, for all \(n \geq N\), it follows that
    \[
      \abs{(s_n + t_n) - (s + t)} = \abs{(s_n - s) + (t_n - t)} \leq \abs{s_n - s} + \abs{t_n - t} < \frac{\varepsilon}{2} + \frac{\varepsilon}{2} = \varepsilon.
    \]
    Therefore, by definition, \(\lim (s_n + t_n) = s + t\).
  \end{subprove}

  \newpage 
  
  \item \textbf{Complete but do not hand in}. Prove that if 
  \((s_n)\) converges to \(s\) and \((t_n)\) converges to 
  \(t\), then \((s_n t_n)\) converges to \(s t\). \\
  \textit{Hint: use the triangle inequality, 
  boundedness of convergent sequences, and addition of \(0 = -s_n t + s_n t\)}.
  \begin{subprove}
    Suppose that \((s_n)\) converges to \(s \in \setR\) and \((t_n)\) converges to \(t \in \setR\) 
    for the sake of direct proof.
    
    Since \((s_n)\) is convergent, by the previous result, \((s_n)\) is bounded.
    Therefore, by \nameref{def:bounded-sequence}, there exists \(M > 0\) such that
    \[
      \abs{s_n} \leq M \quad\text{for all } n \in \setN.
    \]
    
    Let \(\varepsilon > 0\) be given. Since \(M > 0\), it follows that \(\varepsilon / (2M) > 0\) 
    is also an arbitrary positive real number.
    Then, by definition, since \(\lim t_n = t\),
    there exists \(N_1 \in \setN\) such that
    \[
      \abs{t_n - t} < \frac{\varepsilon}{2M} \quad\text{for all } n \geq N_1.
    \]
    
    Since \(\abs{t} \geq 0\), it follows that \(\max\{\abs{t}, 1\} \geq 1 > 0\).
    Since \(\varepsilon > 0\), it follows that \(\varepsilon / (2\max\{\abs{t}, 1\}) > 0\) 
    is also an arbitrary positive real number.
    Then, by definition, since \(\lim s_n = s\),
    there exists \(N_2 \in \setN\) such that
    \[
      \abs{s_n - s} < \frac{\varepsilon}{2\max\{\abs{t}, 1\}} \quad\text{for all } n \geq N_2.
    \]
    
    Let \(N = \max\{N_1, N_2\}\). Then, for all \(n \geq N\), it follows that
    \begin{align*}
      \abs{s_n t_n - s t} &= \abs{s_n t_n - s_n t + s_n t - s t} \\
      &= \abs{s_n(t_n - t) + t(s_n - s)} \\
      &\leq \abs{s_n}\abs{t_n - t} + \abs{t}\abs{s_n - s} \\
      &< M \cdot \frac{\varepsilon}{2M} + \abs{t} \cdot \frac{\varepsilon}{2\max\{\abs{t}, 1\}} \\
      &= \frac{\varepsilon}{2} + \frac{\abs{t}}{\max\{\abs{t}, 1\}} \cdot \frac{\varepsilon}{2} \\
      &\leq \frac{\varepsilon}{2} + \frac{\varepsilon}{2} = \varepsilon.
    \end{align*}
    Therefore, by definition, \(\lim (s_n t_n) = s t\).
  \end{subprove}

\end{enumerate}
